\documentclass[../main.tex]{subfiles}
\begin{document}

\chapter{スレッド設計上の考慮事項}
\lessoninfo{
  video={P2L4},
  textbook={OSTEP \cite{ostep} ch.~6，26--28；Silberschatz ら \cite{silberschatz2018} ch.~4；Kerrisk \cite{kerrisk2010} ch.~20--22，28},
  keywords={ハードプロセス状態，ライトプロセス状態，軽量プロセス，レッドゾーン，適応型ミューテックス，リーパスレッド，割り込みハンドラ，シグナルハンドラ，シグナルマスク，トップハーフ，ボトムハーフ，リアルタイムシグナル，タスク，NPTL},
}

第3章ではカーネルレベルスレッドとユーザレベルスレッド，および一方を他方に
対応付けるモデルを導入し，第4章ではアプリケーションがそれらを利用する
ための POSIX スレッドのインタフェースを扱った．この章では実装の内側を
見る．すなわち，カーネルとユーザレベルのスレッドライブラリがどのような
データ構造を保持するか，2つのレベルがどのように判断を協調させるか，そして
割り込みとシグナルがマルチスレッドとどう関わるかである．説明は SunOS~5.0
のマルチスレッドカーネル \cite{eykholt1992} とそのユーザレベルスレッド
ライブラリ \cite{stein1992} に関する2つの論文に沿って進め，最後に Linux が
スレッドを表現する方法で締めくくる．

\section{スレッドのデータ構造}
\label{sec:l05-structures}

スレッドを支えるには，2つのレベルでデータ構造が必要である．カーネルレベル
スレッドのためのカーネル内のものと，ユーザレベルスレッドのためのユーザ
レベルライブラリ内のものである．\source{P2L4，スレッドのデータ構造：単一
CPU の場合と大規模な場合，ハードプロセス状態とライトプロセス状態，データ
構造を分ける根拠；Eykholt ら \cite{eykholt1992}；OSTEP ch.~26．}まず，すべての
ユーザレベルスレッドを1つのカーネルレベルスレッド上で実行するプロセス
（第3章の多対一モデル）を考える．カーネルはこのプロセスを PCB で記述し，PCB は仮想
アドレスの対応付けと，カーネルレベルスレッドの実行コンテキスト，すなわち
そのレジスタ，スタックポインタ，スタックを保持する．ユーザレベル
ライブラリは，ユーザレベルスレッドごとにデータ構造を保持し，そこにスレッド
の識別子，保存されたレジスタ，スタックを収める．これらのコンテキストの1つを
カーネルレベルスレッドに読み込むことで，ライブラリはどのユーザレベル
スレッドを実行するかを決める．\margin{PCB そのものは第2章で扱った．ここでは
その内容を分割し，1つのプロセスのスレッドがアドレス空間の部分を共有できる
ようにする．}

プロセスが複数のカーネルレベルスレッドを持つ場合，1つの PCB ではすべてを
保持できない．カーネルレベルスレッドはそれぞれ独自の実行コンテキストを
必要とするからである．各スレッドについて PCB 全体を複製すると，プロセスの
すべてのスレッドで同一である仮想アドレスの対応付けまで複製することになる．
そこでカーネルは，カーネルレベルスレッドごとにそのレジスタ，スタック
ポインタ，スタックを収める別個のカーネルレベルスレッドデータ構造を保持し，
PCB にはプロセスに共通の情報を残す．\margin{x86-64 の Linux では，タスクは
それぞれ \SI{16}{\kibi\byte}（4ページ）のカーネルスタックを持つが，
アドレス空間は共有される．カーネル文書 \texttt{x86/kernel-stacks}．}

多数のプロセス，多数のスレッド，複数の CPU がある場合には，これらの構造の
間の関係も記録しなければならない（\cref{fig:l05-structures}）．
\begin{itemize}
  \item ユーザレベルライブラリは，自プロセスのすべてのユーザレベル
    スレッドと，それらを実行できるカーネルレベルスレッドを把握する．
  \item PCB は，そのプロセスのために実行されるカーネルレベルスレッドを
    把握する．
  \item 各カーネルレベルスレッドは自プロセスの PCB を指し，カーネルは
    そのスレッドを実行するときにアドレス空間を見つけられる．
  \item CPU ごとのデータ構造が，とりわけその CPU で現在実行中のスレッドを
    記録し，実行中のカーネルレベルスレッドは自分の CPU を指す．
\end{itemize}

\begin{figure}[htbp]
  \centering
  % Thread-related data structures at scale: user-level thread structures in
% the library, PCBs holding the hard process state, kernel-level threads
% holding the light process state, and CPU structures.
% Usage: % Thread-related data structures at scale: user-level thread structures in
% the library, PCBs holding the hard process state, kernel-level threads
% holding the light process state, and CPU structures.
% Usage: % Thread-related data structures at scale: user-level thread structures in
% the library, PCBs holding the hard process state, kernel-level threads
% holding the light process state, and CPU structures.
% Usage: \input{figures/l05-structures}   (width ~118 mm)
\begin{tikzpicture}[x=1mm, y=1mm, font=\scriptsize\sffamily,
  >={Stealth[length=1.5mm]},
  ult/.style={draw, fill=white, minimum width=10mm, minimum height=5mm,
              inner sep=0.5pt, align=center},
  klt/.style={draw, fill=white, minimum width=16mm, minimum height=11mm,
              inner sep=0.5pt, align=center},
  pcb/.style={draw, fill=ln-box, minimum height=14mm, inner sep=0.5pt,
              align=center},
  cpu/.style={draw, fill=ln-accent!22, minimum width=12mm, minimum height=5mm,
              inner sep=0.5pt},
  lab/.style={font=\scriptsize\sffamily, text=ln-muted},
  map/.style={densely dashed, ln-muted}]
  % user/kernel boundary
  \draw[dashed, ln-rule] (0,36) -- (118,36);
  \node[lab, anchor=south west] at (0,36.5) {\iflangja{ユーザ}{user}};
  \node[lab, anchor=north west] at (0,35.5) {\iflangja{カーネル}{kernel}};
  % ---- process P: library with user-level threads
  \draw[rounded corners=2pt, fill=ln-box] (22,39) rectangle (86,56);
  \node[lab, anchor=north west] at (22.5,55.8)
    {\iflangja{P のユーザレベルスレッドライブラリ}{user-level thread library of P}};
  \foreach \x/\n in {30/1, 44/2, 58/3, 72/4}
    \node[ult] (u\n) at (\x,45) {ULT\,\n};
  % ---- process P: PCB and kernel-level threads
  \node[pcb, minimum width=19mm] (pcbp) at (9.5,24)
    {\iflangja{P の PCB}{PCB of P}\\[1pt]
     \iflangja{ハード状態}{hard state}\\
     \iflangja{（アドレス}{(address}\\
     \iflangja{対応付け）}{mappings)}};
  \foreach \x/\n in {32/1, 54/2, 76/3}
    \node[klt] (k\n) at (\x,24)
      {\textbf{KLT\,\n}\\\iflangja{ライト状態}{light state}\\
       \iflangja{レジスタ等}{regs, stack}};
  \draw[->] (pcbp.east) -- (k1.west);
  \draw[->] (k1.east) -- (k2.west);
  \draw[->] (k2.east) -- (k3.west);
  % user-level threads currently mapped onto kernel-level threads
  \draw[map] (u1.south) -- (k1.north);
  \draw[map] (u3.south) -- (k2.north);
  \draw[map] (u4.south) -- (k3.north);
  % ---- process Q: single-threaded
  \draw[rounded corners=2pt, fill=ln-box] (93,39) rectangle (118,56);
  \node[align=center] at (105.5,47.5)
    {\iflangja{プロセス Q}{process Q}\\\iflangja{（シングル\\スレッド）}{(single-\\threaded)}};
  \node[pcb, minimum width=11mm] (pcbq) at (95.5,24)
    {\iflangja{Q の}{PCB}\\\iflangja{PCB}{of Q}};
  \node[klt, minimum width=14mm] (kq) at (111,24)
    {\textbf{KLT}\\\iflangja{ライト状態}{light state}\\
     \iflangja{レジスタ等}{regs, stack}};
  \draw[->] (pcbq.east) -- (kq.west);
  % ---- CPUs
  \node[cpu] (c0) at (32,5) {CPU\,0};
  \node[cpu] (c1) at (76,5) {CPU\,1};
  \node[cpu] (c2) at (111,5) {CPU\,2};
  \draw[<->] (k1.south) -- (c0.north);
  \draw[<->] (k3.south) -- (c1.north);
  \draw[<->] (kq.south) -- (c2.north);
  \node[lab, anchor=north] at (54,18) {\iflangja{実行中でない}{not running}};
\end{tikzpicture}
   (width ~118 mm)
\begin{tikzpicture}[x=1mm, y=1mm, font=\scriptsize\sffamily,
  >={Stealth[length=1.5mm]},
  ult/.style={draw, fill=white, minimum width=10mm, minimum height=5mm,
              inner sep=0.5pt, align=center},
  klt/.style={draw, fill=white, minimum width=16mm, minimum height=11mm,
              inner sep=0.5pt, align=center},
  pcb/.style={draw, fill=ln-box, minimum height=14mm, inner sep=0.5pt,
              align=center},
  cpu/.style={draw, fill=ln-accent!22, minimum width=12mm, minimum height=5mm,
              inner sep=0.5pt},
  lab/.style={font=\scriptsize\sffamily, text=ln-muted},
  map/.style={densely dashed, ln-muted}]
  % user/kernel boundary
  \draw[dashed, ln-rule] (0,36) -- (118,36);
  \node[lab, anchor=south west] at (0,36.5) {\iflangja{ユーザ}{user}};
  \node[lab, anchor=north west] at (0,35.5) {\iflangja{カーネル}{kernel}};
  % ---- process P: library with user-level threads
  \draw[rounded corners=2pt, fill=ln-box] (22,39) rectangle (86,56);
  \node[lab, anchor=north west] at (22.5,55.8)
    {\iflangja{P のユーザレベルスレッドライブラリ}{user-level thread library of P}};
  \foreach \x/\n in {30/1, 44/2, 58/3, 72/4}
    \node[ult] (u\n) at (\x,45) {ULT\,\n};
  % ---- process P: PCB and kernel-level threads
  \node[pcb, minimum width=19mm] (pcbp) at (9.5,24)
    {\iflangja{P の PCB}{PCB of P}\\[1pt]
     \iflangja{ハード状態}{hard state}\\
     \iflangja{（アドレス}{(address}\\
     \iflangja{対応付け）}{mappings)}};
  \foreach \x/\n in {32/1, 54/2, 76/3}
    \node[klt] (k\n) at (\x,24)
      {\textbf{KLT\,\n}\\\iflangja{ライト状態}{light state}\\
       \iflangja{レジスタ等}{regs, stack}};
  \draw[->] (pcbp.east) -- (k1.west);
  \draw[->] (k1.east) -- (k2.west);
  \draw[->] (k2.east) -- (k3.west);
  % user-level threads currently mapped onto kernel-level threads
  \draw[map] (u1.south) -- (k1.north);
  \draw[map] (u3.south) -- (k2.north);
  \draw[map] (u4.south) -- (k3.north);
  % ---- process Q: single-threaded
  \draw[rounded corners=2pt, fill=ln-box] (93,39) rectangle (118,56);
  \node[align=center] at (105.5,47.5)
    {\iflangja{プロセス Q}{process Q}\\\iflangja{（シングル\\スレッド）}{(single-\\threaded)}};
  \node[pcb, minimum width=11mm] (pcbq) at (95.5,24)
    {\iflangja{Q の}{PCB}\\\iflangja{PCB}{of Q}};
  \node[klt, minimum width=14mm] (kq) at (111,24)
    {\textbf{KLT}\\\iflangja{ライト状態}{light state}\\
     \iflangja{レジスタ等}{regs, stack}};
  \draw[->] (pcbq.east) -- (kq.west);
  % ---- CPUs
  \node[cpu] (c0) at (32,5) {CPU\,0};
  \node[cpu] (c1) at (76,5) {CPU\,1};
  \node[cpu] (c2) at (111,5) {CPU\,2};
  \draw[<->] (k1.south) -- (c0.north);
  \draw[<->] (k3.south) -- (c1.north);
  \draw[<->] (kq.south) -- (c2.north);
  \node[lab, anchor=north] at (54,18) {\iflangja{実行中でない}{not running}};
\end{tikzpicture}
   (width ~118 mm)
\begin{tikzpicture}[x=1mm, y=1mm, font=\scriptsize\sffamily,
  >={Stealth[length=1.5mm]},
  ult/.style={draw, fill=white, minimum width=10mm, minimum height=5mm,
              inner sep=0.5pt, align=center},
  klt/.style={draw, fill=white, minimum width=16mm, minimum height=11mm,
              inner sep=0.5pt, align=center},
  pcb/.style={draw, fill=ln-box, minimum height=14mm, inner sep=0.5pt,
              align=center},
  cpu/.style={draw, fill=ln-accent!22, minimum width=12mm, minimum height=5mm,
              inner sep=0.5pt},
  lab/.style={font=\scriptsize\sffamily, text=ln-muted},
  map/.style={densely dashed, ln-muted}]
  % user/kernel boundary
  \draw[dashed, ln-rule] (0,36) -- (118,36);
  \node[lab, anchor=south west] at (0,36.5) {\iflangja{ユーザ}{user}};
  \node[lab, anchor=north west] at (0,35.5) {\iflangja{カーネル}{kernel}};
  % ---- process P: library with user-level threads
  \draw[rounded corners=2pt, fill=ln-box] (22,39) rectangle (86,56);
  \node[lab, anchor=north west] at (22.5,55.8)
    {\iflangja{P のユーザレベルスレッドライブラリ}{user-level thread library of P}};
  \foreach \x/\n in {30/1, 44/2, 58/3, 72/4}
    \node[ult] (u\n) at (\x,45) {ULT\,\n};
  % ---- process P: PCB and kernel-level threads
  \node[pcb, minimum width=19mm] (pcbp) at (9.5,24)
    {\iflangja{P の PCB}{PCB of P}\\[1pt]
     \iflangja{ハード状態}{hard state}\\
     \iflangja{（アドレス}{(address}\\
     \iflangja{対応付け）}{mappings)}};
  \foreach \x/\n in {32/1, 54/2, 76/3}
    \node[klt] (k\n) at (\x,24)
      {\textbf{KLT\,\n}\\\iflangja{ライト状態}{light state}\\
       \iflangja{レジスタ等}{regs, stack}};
  \draw[->] (pcbp.east) -- (k1.west);
  \draw[->] (k1.east) -- (k2.west);
  \draw[->] (k2.east) -- (k3.west);
  % user-level threads currently mapped onto kernel-level threads
  \draw[map] (u1.south) -- (k1.north);
  \draw[map] (u3.south) -- (k2.north);
  \draw[map] (u4.south) -- (k3.north);
  % ---- process Q: single-threaded
  \draw[rounded corners=2pt, fill=ln-box] (93,39) rectangle (118,56);
  \node[align=center] at (105.5,47.5)
    {\iflangja{プロセス Q}{process Q}\\\iflangja{（シングル\\スレッド）}{(single-\\threaded)}};
  \node[pcb, minimum width=11mm] (pcbq) at (95.5,24)
    {\iflangja{Q の}{PCB}\\\iflangja{PCB}{of Q}};
  \node[klt, minimum width=14mm] (kq) at (111,24)
    {\textbf{KLT}\\\iflangja{ライト状態}{light state}\\
     \iflangja{レジスタ等}{regs, stack}};
  \draw[->] (pcbq.east) -- (kq.west);
  % ---- CPUs
  \node[cpu] (c0) at (32,5) {CPU\,0};
  \node[cpu] (c1) at (76,5) {CPU\,1};
  \node[cpu] (c2) at (111,5) {CPU\,2};
  \draw[<->] (k1.south) -- (c0.north);
  \draw[<->] (k3.south) -- (c1.north);
  \draw[<->] (kq.south) -- (c2.north);
  \node[lab, anchor=north] at (54,18) {\iflangja{実行中でない}{not running}};
\end{tikzpicture}

  \caption{大規模な場合のスレッドのデータ構造．プロセスの PCB は，その
    すべてのスレッドが共有する状態と，カーネルレベルスレッド（KLT）の
    リストを保持する．各 KLT は独自の状態を保持し，実行中は CPU に
    結び付けられる．破線は，ライブラリが現在どのユーザレベルスレッド
    （ULT）をどの KLT に配置しているかを示す．KLT から PCB への逆向きの
    ポインタは省いてある．}
  \label{fig:l05-structures}
\end{figure}

\subsection{ハードプロセス状態とライトプロセス状態}

カーネルが同じプロセスの2つのカーネルレベルスレッドを切り替えるとき，仮想
アドレスの対応付けは有効なままであり，保存・復元しなければならないのは
個々のスレッドに属する状態だけである．\tip{プロセスのすべてのスレッドに
ついて有効ならハード状態，現在実行中のスレッドについてだけ有効ならライト
状態である．}その状態の一部はユーザレベル
スレッドに関わる．どのシグナルがマスクされているか，および実行中の
システムコールの引数は，プロセス全体ではなく，そのカーネルレベルスレッド上
で現在実行しているユーザレベルスレッドに属する．こうしてプロセスの状態は
2つの部分に分かれる．

\needspace{9\baselineskip}
\begin{definition}[ハードプロセス状態とライトプロセス状態]
  \term[はーどぷろせすじょうたい]{ハードプロセス状態}[hard process state]%
  とは，プロセスの状態のうち，その中で実行されるすべてのユーザレベル
  スレッドに関係する部分であり，仮想アドレスの対応付けなどがこれにあたる．
  \term[らいとぷろせすじょうたい]{ライトプロセス状態}[light process state]%
  とは，特定のカーネルレベルスレッドに現在結び付けられているユーザレベル
  スレッドの部分集合にのみ関係する部分であり，シグナルマスクやシステム
  コールの引数などがこれにあたる．\caution{この2つの用語は Solaris の
    文献と講義のものである．教科書では通常，プロセスごとの状態，スレッド
    ごとの状態と呼ぶ．}
\end{definition}

\subsection{データ構造を複数に分ける理由}

1つの PCB と，相互に参照する小さな構造の集合とでは，コストが異なる
（\cref{tab:l05-rationale}）．1つの PCB は大きな連続した構造であり，その
内容が共有されている場合でも実行コンテキストごとに専有される．コンテキスト
スイッチのたびに全体が保存・復元され，状態が変化すればそのたびに更新される．
ポインタで互いを参照する小さな構造であれば，共有される状態は1つだけ保持
すればよく，コンテキストスイッチで保存・復元するのは変化する部分だけで済み，
ユーザレベルライブラリは自ら管理する部分をカーネルを介さずに更新できる．
この分割はスケーラビリティ，性能，柔軟性を改善し，オーバーヘッドを減らす．

\begin{table}[htbp]
  \centering
  \caption{1つの PCB と複数の小さなデータ構造の比較．}
  \label{tab:l05-rationale}
  \small
  \begin{tabular}{@{}>{\raggedright\arraybackslash}p{0.19\linewidth}>{\raggedright\arraybackslash}p{0.37\linewidth}>{\raggedright\arraybackslash}p{0.37\linewidth}@{}}
    \toprule
    & 1つの PCB & 複数のデータ構造 \\
    \midrule
    構造 & 大きな連続した1つの構造
         & ポインタで連結された小さな構造 \\
    \addlinespace
    共有 & 実行コンテキストごとに専有
         & 共有される状態は1つだけ保持 \\
    \addlinespace
    コンテキスト\newline スイッチ & 全体を保存・復元
         & 変化する部分だけ \\
    \addlinespace
    更新 & 変更のたびに PCB を書き換え
         & ライブラリは自分の部分だけを更新 \\
    \midrule
    帰結 & スケーラビリティが低く，オーバーヘッドが大きく，性能が低く，
           柔軟性に乏しい
         & スケーラビリティが高く，オーバーヘッドが小さく，性能が高く，
           柔軟性に富む \\
    \bottomrule
  \end{tabular}
\end{table}

\section{SunOS~5.0 のスレッド}
\label{sec:l05-solaris}

SunOS~5.0，すなわち Solaris~2 のカーネルは共有メモリ型マルチプロセッサ
向けに設計されており，カーネル自体がマルチスレッドである．\source{P2L4，
SunOS~5.0 のスレッドモデル，ユーザレベルスレッドのデータ構造，カーネル
レベルのデータ構造；Eykholt ら \cite{eykholt1992}；Stein と Shah
\cite{stein1992}．}そのスレッドモデル（\cref{fig:l05-sunos}）は次の要素
からなる．
\begin{itemize}
  \item プロセスはシングルスレッドでもマルチスレッドでもよい．ユーザ
    レベルライブラリは，バウンドスレッドを含む多対多モデルと一対一モデル
    の両方に対応し，プロセスはすべてのユーザレベルスレッドを1つの軽量
    プロセス上で実行することもできる．
  \item ユーザレベルスレッドを実行するカーネルレベルスレッドは，すべて
    軽量プロセスと結び付けられる．
  \item それ以外のカーネルレベルスレッドは軽量プロセスを持たず，カーネル
    自身の仕事を行う．
  \item カーネルレベルスケジューラが，カーネルレベルスレッドを CPU に
    スケジュールする．
\end{itemize}

\begin{definition}[軽量プロセス]
  SunOS~5.0 における\term[けいりょうぷろせす]{軽量プロセス}[lightweight process]%
  \index{LWP|see{軽量プロセス}}（LWP）とは，ユーザレベルスレッドを実行する
  カーネルレベルスレッドに結び付けられたカーネルのデータ構造であり，その
  実行コンテキストのライトプロセス状態を保持する．ユーザレベルライブラリ
  から見ると，LWP はユーザレベルスレッドをスケジュールする先の仮想 CPU を
  表す．
\end{definition}

\begin{figure}[htbp]
  \centering
  % The SunOS 5.0 threading model: user-level threads, lightweight processes,
% kernel-level threads, the kernel-level scheduler and the CPUs.
% Usage: % The SunOS 5.0 threading model: user-level threads, lightweight processes,
% kernel-level threads, the kernel-level scheduler and the CPUs.
% Usage: % The SunOS 5.0 threading model: user-level threads, lightweight processes,
% kernel-level threads, the kernel-level scheduler and the CPUs.
% Usage: \input{figures/l05-sunos}   (width ~118 mm)
\begin{tikzpicture}[x=1mm, y=1mm, font=\scriptsize\sffamily,
  ut/.style={draw, circle, fill=white, minimum size=3.6mm, inner sep=0pt},
  lwp/.style={draw, rounded corners=1pt, fill=ln-accent!15, minimum width=6.5mm,
              minimum height=3.6mm, inner sep=0pt, font=\tiny\sffamily},
  kt/.style={draw, rectangle, fill=ln-box, minimum size=3.6mm, inner sep=0pt},
  cpu/.style={draw, fill=ln-accent!22, minimum width=12mm, minimum height=4.5mm,
              inner sep=0pt},
  proc/.style={draw, dashed, rounded corners=2pt, ln-muted},
  lab/.style={font=\scriptsize\sffamily, text=ln-muted},
  ttl/.style={font=\scriptsize\sffamily\bfseries, anchor=base}]
  % user/kernel boundary runs through the LWPs
  \draw[dashed, ln-rule] (0,22) -- (118,22);
  \node[lab, anchor=south west] at (0,22.5) {\iflangja{ユーザ}{user}};
  \node[lab, anchor=north west] at (0,21.5) {\iflangja{カーネル}{kernel}};
  % ---- P1: many-to-many running on a single LWP
  \draw[proc] (12,17) rectangle (36,39);
  \node[ttl] at (24,40.5) {\iflangja{P1：LWP 1 個}{P1: one LWP}};
  \node[lwp] (l1) at (24,22) {LWP};
  \foreach \x in {17,24,31} { \node[ut] (u) at (\x,33) {}; \draw (u) -- (l1.north); }
  \node[kt] (k1) at (24,12) {};
  \draw (l1) -- (k1);
  % ---- P2: many-to-many with a bound thread
  \draw[proc] (40,17) rectangle (82,39);
  \node[ttl] at (61,40.5) {\iflangja{P2：多対多}{P2: many-to-many}};
  \node[lwp] (l2a) at (49,22) {LWP};
  \node[lwp] (l2b) at (60,22) {LWP};
  \node[ut] (u2a) at (45,33) {};
  \node[ut] (u2b) at (54.5,33) {};
  \node[ut] (u2c) at (64,33) {};
  \draw (u2a) -- (l2a.north);
  \draw (u2b) -- (l2a.north);
  \draw (u2b) -- (l2b.north);
  \draw (u2c) -- (l2b.north);
  \node[kt] (k2a) at (49,12) {};
  \node[kt] (k2b) at (60,12) {};
  \draw (l2a) -- (k2a);
  \draw (l2b) -- (k2b);
  \node[ut, thick, draw=ln-accent] (ub) at (75,33) {};
  \node[lwp, thick, draw=ln-accent] (lb) at (75,22) {LWP};
  \node[kt, thick, draw=ln-accent] (kb) at (75,12) {};
  \draw[very thick, ln-accent] (ub) -- (lb) -- (kb);
  \node[lab, text=ln-accent, anchor=south] at (75,35) {\iflangja{バウンド}{bound}};
  % ---- P3: one-to-one
  \draw[proc] (86,17) rectangle (106,39);
  \node[ttl] at (96,40.5) {\iflangja{P3：一対一}{P3: one-to-one}};
  \foreach \x/\n in {91/a, 101/b} {
    \node[ut] (u) at (\x,33) {};
    \node[lwp] (l) at (\x,22) {LWP};
    \node[kt] (k3\n) at (\x,12) {};
    \draw (u) -- (l) -- (k3\n);
  }
  % ---- kernel threads without LWP
  \node[kt] (kk1) at (111,12) {};
  \node[kt] (kk2) at (116,12) {};
  % ---- scheduler and CPUs
  \draw[rounded corners=1pt, fill=ln-box] (12,2) rectangle (118,7);
  \node at (65,4.5) {\iflangja{カーネルレベルスケジューラ}{kernel-level scheduler}};
  \foreach \k in {k1, k2a, k2b, kb, k3a, k3b, kk1, kk2}
    \draw (\k.south) -- (\k.south |- 0,7);
  \foreach \x/\n in {35/0, 65/1, 95/2} {
    \node[cpu] (c) at (\x,-6) {CPU\,\n};
    \draw (c.north) -- (\x,2);
  }
  % ---- legend
  \node[ut] at (14,-15) {};
  \node[lab, anchor=west] at (16.5,-15) {\iflangja{ユーザレベルスレッド}{user-level thread}};
  \node[lwp] at (48,-15) {LWP};
  \node[lab, anchor=west] at (52,-15) {\iflangja{軽量プロセス}{lightweight process}};
  \node[kt] at (84,-15) {};
  \node[lab, anchor=west] at (86.5,-15) {\iflangja{カーネルレベルスレッド}{kernel-level thread}};
\end{tikzpicture}
   (width ~118 mm)
\begin{tikzpicture}[x=1mm, y=1mm, font=\scriptsize\sffamily,
  ut/.style={draw, circle, fill=white, minimum size=3.6mm, inner sep=0pt},
  lwp/.style={draw, rounded corners=1pt, fill=ln-accent!15, minimum width=6.5mm,
              minimum height=3.6mm, inner sep=0pt, font=\tiny\sffamily},
  kt/.style={draw, rectangle, fill=ln-box, minimum size=3.6mm, inner sep=0pt},
  cpu/.style={draw, fill=ln-accent!22, minimum width=12mm, minimum height=4.5mm,
              inner sep=0pt},
  proc/.style={draw, dashed, rounded corners=2pt, ln-muted},
  lab/.style={font=\scriptsize\sffamily, text=ln-muted},
  ttl/.style={font=\scriptsize\sffamily\bfseries, anchor=base}]
  % user/kernel boundary runs through the LWPs
  \draw[dashed, ln-rule] (0,22) -- (118,22);
  \node[lab, anchor=south west] at (0,22.5) {\iflangja{ユーザ}{user}};
  \node[lab, anchor=north west] at (0,21.5) {\iflangja{カーネル}{kernel}};
  % ---- P1: many-to-many running on a single LWP
  \draw[proc] (12,17) rectangle (36,39);
  \node[ttl] at (24,40.5) {\iflangja{P1：LWP 1 個}{P1: one LWP}};
  \node[lwp] (l1) at (24,22) {LWP};
  \foreach \x in {17,24,31} { \node[ut] (u) at (\x,33) {}; \draw (u) -- (l1.north); }
  \node[kt] (k1) at (24,12) {};
  \draw (l1) -- (k1);
  % ---- P2: many-to-many with a bound thread
  \draw[proc] (40,17) rectangle (82,39);
  \node[ttl] at (61,40.5) {\iflangja{P2：多対多}{P2: many-to-many}};
  \node[lwp] (l2a) at (49,22) {LWP};
  \node[lwp] (l2b) at (60,22) {LWP};
  \node[ut] (u2a) at (45,33) {};
  \node[ut] (u2b) at (54.5,33) {};
  \node[ut] (u2c) at (64,33) {};
  \draw (u2a) -- (l2a.north);
  \draw (u2b) -- (l2a.north);
  \draw (u2b) -- (l2b.north);
  \draw (u2c) -- (l2b.north);
  \node[kt] (k2a) at (49,12) {};
  \node[kt] (k2b) at (60,12) {};
  \draw (l2a) -- (k2a);
  \draw (l2b) -- (k2b);
  \node[ut, thick, draw=ln-accent] (ub) at (75,33) {};
  \node[lwp, thick, draw=ln-accent] (lb) at (75,22) {LWP};
  \node[kt, thick, draw=ln-accent] (kb) at (75,12) {};
  \draw[very thick, ln-accent] (ub) -- (lb) -- (kb);
  \node[lab, text=ln-accent, anchor=south] at (75,35) {\iflangja{バウンド}{bound}};
  % ---- P3: one-to-one
  \draw[proc] (86,17) rectangle (106,39);
  \node[ttl] at (96,40.5) {\iflangja{P3：一対一}{P3: one-to-one}};
  \foreach \x/\n in {91/a, 101/b} {
    \node[ut] (u) at (\x,33) {};
    \node[lwp] (l) at (\x,22) {LWP};
    \node[kt] (k3\n) at (\x,12) {};
    \draw (u) -- (l) -- (k3\n);
  }
  % ---- kernel threads without LWP
  \node[kt] (kk1) at (111,12) {};
  \node[kt] (kk2) at (116,12) {};
  % ---- scheduler and CPUs
  \draw[rounded corners=1pt, fill=ln-box] (12,2) rectangle (118,7);
  \node at (65,4.5) {\iflangja{カーネルレベルスケジューラ}{kernel-level scheduler}};
  \foreach \k in {k1, k2a, k2b, kb, k3a, k3b, kk1, kk2}
    \draw (\k.south) -- (\k.south |- 0,7);
  \foreach \x/\n in {35/0, 65/1, 95/2} {
    \node[cpu] (c) at (\x,-6) {CPU\,\n};
    \draw (c.north) -- (\x,2);
  }
  % ---- legend
  \node[ut] at (14,-15) {};
  \node[lab, anchor=west] at (16.5,-15) {\iflangja{ユーザレベルスレッド}{user-level thread}};
  \node[lwp] at (48,-15) {LWP};
  \node[lab, anchor=west] at (52,-15) {\iflangja{軽量プロセス}{lightweight process}};
  \node[kt] at (84,-15) {};
  \node[lab, anchor=west] at (86.5,-15) {\iflangja{カーネルレベルスレッド}{kernel-level thread}};
\end{tikzpicture}
   (width ~118 mm)
\begin{tikzpicture}[x=1mm, y=1mm, font=\scriptsize\sffamily,
  ut/.style={draw, circle, fill=white, minimum size=3.6mm, inner sep=0pt},
  lwp/.style={draw, rounded corners=1pt, fill=ln-accent!15, minimum width=6.5mm,
              minimum height=3.6mm, inner sep=0pt, font=\tiny\sffamily},
  kt/.style={draw, rectangle, fill=ln-box, minimum size=3.6mm, inner sep=0pt},
  cpu/.style={draw, fill=ln-accent!22, minimum width=12mm, minimum height=4.5mm,
              inner sep=0pt},
  proc/.style={draw, dashed, rounded corners=2pt, ln-muted},
  lab/.style={font=\scriptsize\sffamily, text=ln-muted},
  ttl/.style={font=\scriptsize\sffamily\bfseries, anchor=base}]
  % user/kernel boundary runs through the LWPs
  \draw[dashed, ln-rule] (0,22) -- (118,22);
  \node[lab, anchor=south west] at (0,22.5) {\iflangja{ユーザ}{user}};
  \node[lab, anchor=north west] at (0,21.5) {\iflangja{カーネル}{kernel}};
  % ---- P1: many-to-many running on a single LWP
  \draw[proc] (12,17) rectangle (36,39);
  \node[ttl] at (24,40.5) {\iflangja{P1：LWP 1 個}{P1: one LWP}};
  \node[lwp] (l1) at (24,22) {LWP};
  \foreach \x in {17,24,31} { \node[ut] (u) at (\x,33) {}; \draw (u) -- (l1.north); }
  \node[kt] (k1) at (24,12) {};
  \draw (l1) -- (k1);
  % ---- P2: many-to-many with a bound thread
  \draw[proc] (40,17) rectangle (82,39);
  \node[ttl] at (61,40.5) {\iflangja{P2：多対多}{P2: many-to-many}};
  \node[lwp] (l2a) at (49,22) {LWP};
  \node[lwp] (l2b) at (60,22) {LWP};
  \node[ut] (u2a) at (45,33) {};
  \node[ut] (u2b) at (54.5,33) {};
  \node[ut] (u2c) at (64,33) {};
  \draw (u2a) -- (l2a.north);
  \draw (u2b) -- (l2a.north);
  \draw (u2b) -- (l2b.north);
  \draw (u2c) -- (l2b.north);
  \node[kt] (k2a) at (49,12) {};
  \node[kt] (k2b) at (60,12) {};
  \draw (l2a) -- (k2a);
  \draw (l2b) -- (k2b);
  \node[ut, thick, draw=ln-accent] (ub) at (75,33) {};
  \node[lwp, thick, draw=ln-accent] (lb) at (75,22) {LWP};
  \node[kt, thick, draw=ln-accent] (kb) at (75,12) {};
  \draw[very thick, ln-accent] (ub) -- (lb) -- (kb);
  \node[lab, text=ln-accent, anchor=south] at (75,35) {\iflangja{バウンド}{bound}};
  % ---- P3: one-to-one
  \draw[proc] (86,17) rectangle (106,39);
  \node[ttl] at (96,40.5) {\iflangja{P3：一対一}{P3: one-to-one}};
  \foreach \x/\n in {91/a, 101/b} {
    \node[ut] (u) at (\x,33) {};
    \node[lwp] (l) at (\x,22) {LWP};
    \node[kt] (k3\n) at (\x,12) {};
    \draw (u) -- (l) -- (k3\n);
  }
  % ---- kernel threads without LWP
  \node[kt] (kk1) at (111,12) {};
  \node[kt] (kk2) at (116,12) {};
  % ---- scheduler and CPUs
  \draw[rounded corners=1pt, fill=ln-box] (12,2) rectangle (118,7);
  \node at (65,4.5) {\iflangja{カーネルレベルスケジューラ}{kernel-level scheduler}};
  \foreach \k in {k1, k2a, k2b, kb, k3a, k3b, kk1, kk2}
    \draw (\k.south) -- (\k.south |- 0,7);
  \foreach \x/\n in {35/0, 65/1, 95/2} {
    \node[cpu] (c) at (\x,-6) {CPU\,\n};
    \draw (c.north) -- (\x,2);
  }
  % ---- legend
  \node[ut] at (14,-15) {};
  \node[lab, anchor=west] at (16.5,-15) {\iflangja{ユーザレベルスレッド}{user-level thread}};
  \node[lwp] at (48,-15) {LWP};
  \node[lab, anchor=west] at (52,-15) {\iflangja{軽量プロセス}{lightweight process}};
  \node[kt] at (84,-15) {};
  \node[lab, anchor=west] at (86.5,-15) {\iflangja{カーネルレベルスレッド}{kernel-level thread}};
\end{tikzpicture}

  \caption{SunOS~5.0 のスレッドモデル．ユーザレベルスレッドを実行する
    カーネルレベルスレッドはそれぞれ LWP を持つ．右側の2つのカーネル
    レベルスレッドはカーネルの仕事を行い，LWP を持たない．P1 はすべての
    ユーザレベルスレッドを1つの LWP 上で実行している．カーネルレベル
    スケジューラがカーネルレベルスレッドを CPU に配置する．}
  \label{fig:l05-sunos}
\end{figure}

\subsection{ユーザレベルスレッドのデータ構造}

Stein と Shah のライブラリは，スレッドを生成するとスレッド識別子を返す．
この識別子はスレッドデータ構造へのポインタではなく，構造へのポインタの表
への添字である．この間接参照は古くなった識別子から保護する．スレッドが
すでに終了してその構造のメモリが再利用されている場合，直接のポインタは
無関係なデータを黙って指してしまうが，表のエントリにはスレッドの状態を
記録できるので，ライブラリは意味のあるエラーを返せる．

スレッドデータ構造は，実行コンテキスト（スタックポインタを含むレジスタ），
シグナルマスク，優先度，\emph{スレッドローカルストレージ}，スタックを含む．
スレッドローカルストレージは，スレッド関数がスレッドごとに固有と宣言した
変数を保持する．これらはコンパイル時に分かっているので，コンパイラが
すべてのスレッドに場所を確保できる．スタックのサイズは，ライブラリの既定値
によって，あるいはアプリケーションによって定まる．すべての部分のサイズが
スレッドの生成時に分かっているので，ライブラリは連続したメモリ上に次々と
スレッドの構造を隣り合わせに割り当てることができ，局所性がよくなる．

連続割り当てには危険がある．ライブラリはスタックがどこまで伸びるかを制御
できず，オペレーティングシステムはそもそもユーザレベルスレッドを知らない．
したがって，スタックがあふれたスレッドは隣のスレッドのデータ構造に書き
込んでしまい，その被害は，その別のスレッドが実行されて異常を起こすとき
に---原因から遠く離れたところで---初めて現れる．

\begin{definition}[レッドゾーン]
  \term[れっどぞーん]{レッドゾーン}[red zone]とは，あるスレッドのスタックと
  隣のスレッドのデータ構造の間にある，メモリに対応付けられていない仮想
  アドレス空間の領域であり，そこへのアクセスはすべてフォールトを
  起こす．\caution{同じ語の別の意味：x86-64 の ABI における\emph{レッド
    ゾーン}は，スタックポインタの下の \SI{128}{\byte} であり，リーフ関数が
    スタックポインタを動かさずに使ってよい領域を指す．}
\end{definition}

スタックがレッドゾーンに伸びると直ちに，あふれさせた当のスレッドでフォールト
が起こるので，誤りの診断が容易になる（\cref{fig:l05-ult-layout}）．

\begin{figure}[htbp]
  \centering
  % User-level thread data structures allocated contiguously, with a red zone
% below each stack, and the table of pointers indexed by thread identifier.
% Usage: % User-level thread data structures allocated contiguously, with a red zone
% below each stack, and the table of pointers indexed by thread identifier.
% Usage: % User-level thread data structures allocated contiguously, with a red zone
% below each stack, and the table of pointers indexed by thread identifier.
% Usage: \input{figures/l05-ult-layout}   (width ~116 mm)
\begin{tikzpicture}[x=1mm, y=1mm, font=\scriptsize\sffamily,
  >={Stealth[length=1.5mm]},
  lab/.style={font=\scriptsize\sffamily, text=ln-muted, align=center},
  cell/.style={draw, fill=white, minimum width=8mm, minimum height=5mm,
               inner sep=0pt}]
  \node[lab] at (3,4.5) {\ldots};
  \foreach \o/\t in {8/1, 60/2} {
    \draw[fill=ln-caution!25] (\o,0) rectangle ++(5,9);
    \draw[fill=white] (\o+5,0) rectangle ++(30,9);
    \node at (\o+20,6.2) {\iflangja{T\t のスタック}{stack of T\t}};
    \draw[->, ln-accent, thick] (\o+29,2.5) -- (\o+11,2.5);
    \draw[fill=ln-box] (\o+35,0) rectangle ++(6,9);
    \node at (\o+38,4.5) {TLS};
    \draw[fill=ln-accent!22] (\o+41,0) rectangle ++(11,9);
    \node[align=center] at (\o+46.5,4.5) {T\t\\\iflangja{構造体}{struct}};
    \node[lab, text=ln-caution, anchor=north] at (\o+2.5,-0.8)
      {\iflangja{レッドゾーン}{red zone}};
  }
  \node[lab] at (115,4.5) {\ldots};
  % address axis
  \draw[->, ln-muted] (8,-7.5) -- (112,-7.5);
  \node[lab, anchor=north west] at (8,-8) {\iflangja{下位アドレス}{lower addresses}};
  \node[lab, anchor=north east] at (112,-8) {\iflangja{上位アドレス}{higher addresses}};
  \node[lab, anchor=north] at (60,-8)
    {\iflangja{スタックは下位アドレスへ伸びる}{stacks grow toward lower addresses}};
  % table of pointers indexed by thread id
  \node[lab, anchor=east] at (33,24) {\iflangja{ポインタの表}{table of pointers}};
  \foreach \i/\x in {1/38, 2/46, 3/54} {
    \node[cell] (t\i) at (\x,24) {};
    \node[lab, anchor=south] at (\x,26.5) {\i};
  }
  \node[cell] at (62,24) {\ldots};
  \node[lab, anchor=west] at (67,24) {\iflangja{（添字 = スレッド ID）}{(index = thread ID)}};
  \fill (t1) circle (0.6);
  \fill (t2) circle (0.6);
  \draw[->] (t1.center) -- (54.5,9);
  \draw[->] (t2.center) .. controls (46,15) and (100,16) .. (106.5,9);
\end{tikzpicture}
   (width ~116 mm)
\begin{tikzpicture}[x=1mm, y=1mm, font=\scriptsize\sffamily,
  >={Stealth[length=1.5mm]},
  lab/.style={font=\scriptsize\sffamily, text=ln-muted, align=center},
  cell/.style={draw, fill=white, minimum width=8mm, minimum height=5mm,
               inner sep=0pt}]
  \node[lab] at (3,4.5) {\ldots};
  \foreach \o/\t in {8/1, 60/2} {
    \draw[fill=ln-caution!25] (\o,0) rectangle ++(5,9);
    \draw[fill=white] (\o+5,0) rectangle ++(30,9);
    \node at (\o+20,6.2) {\iflangja{T\t のスタック}{stack of T\t}};
    \draw[->, ln-accent, thick] (\o+29,2.5) -- (\o+11,2.5);
    \draw[fill=ln-box] (\o+35,0) rectangle ++(6,9);
    \node at (\o+38,4.5) {TLS};
    \draw[fill=ln-accent!22] (\o+41,0) rectangle ++(11,9);
    \node[align=center] at (\o+46.5,4.5) {T\t\\\iflangja{構造体}{struct}};
    \node[lab, text=ln-caution, anchor=north] at (\o+2.5,-0.8)
      {\iflangja{レッドゾーン}{red zone}};
  }
  \node[lab] at (115,4.5) {\ldots};
  % address axis
  \draw[->, ln-muted] (8,-7.5) -- (112,-7.5);
  \node[lab, anchor=north west] at (8,-8) {\iflangja{下位アドレス}{lower addresses}};
  \node[lab, anchor=north east] at (112,-8) {\iflangja{上位アドレス}{higher addresses}};
  \node[lab, anchor=north] at (60,-8)
    {\iflangja{スタックは下位アドレスへ伸びる}{stacks grow toward lower addresses}};
  % table of pointers indexed by thread id
  \node[lab, anchor=east] at (33,24) {\iflangja{ポインタの表}{table of pointers}};
  \foreach \i/\x in {1/38, 2/46, 3/54} {
    \node[cell] (t\i) at (\x,24) {};
    \node[lab, anchor=south] at (\x,26.5) {\i};
  }
  \node[cell] at (62,24) {\ldots};
  \node[lab, anchor=west] at (67,24) {\iflangja{（添字 = スレッド ID）}{(index = thread ID)}};
  \fill (t1) circle (0.6);
  \fill (t2) circle (0.6);
  \draw[->] (t1.center) -- (54.5,9);
  \draw[->] (t2.center) .. controls (46,15) and (100,16) .. (106.5,9);
\end{tikzpicture}
   (width ~116 mm)
\begin{tikzpicture}[x=1mm, y=1mm, font=\scriptsize\sffamily,
  >={Stealth[length=1.5mm]},
  lab/.style={font=\scriptsize\sffamily, text=ln-muted, align=center},
  cell/.style={draw, fill=white, minimum width=8mm, minimum height=5mm,
               inner sep=0pt}]
  \node[lab] at (3,4.5) {\ldots};
  \foreach \o/\t in {8/1, 60/2} {
    \draw[fill=ln-caution!25] (\o,0) rectangle ++(5,9);
    \draw[fill=white] (\o+5,0) rectangle ++(30,9);
    \node at (\o+20,6.2) {\iflangja{T\t のスタック}{stack of T\t}};
    \draw[->, ln-accent, thick] (\o+29,2.5) -- (\o+11,2.5);
    \draw[fill=ln-box] (\o+35,0) rectangle ++(6,9);
    \node at (\o+38,4.5) {TLS};
    \draw[fill=ln-accent!22] (\o+41,0) rectangle ++(11,9);
    \node[align=center] at (\o+46.5,4.5) {T\t\\\iflangja{構造体}{struct}};
    \node[lab, text=ln-caution, anchor=north] at (\o+2.5,-0.8)
      {\iflangja{レッドゾーン}{red zone}};
  }
  \node[lab] at (115,4.5) {\ldots};
  % address axis
  \draw[->, ln-muted] (8,-7.5) -- (112,-7.5);
  \node[lab, anchor=north west] at (8,-8) {\iflangja{下位アドレス}{lower addresses}};
  \node[lab, anchor=north east] at (112,-8) {\iflangja{上位アドレス}{higher addresses}};
  \node[lab, anchor=north] at (60,-8)
    {\iflangja{スタックは下位アドレスへ伸びる}{stacks grow toward lower addresses}};
  % table of pointers indexed by thread id
  \node[lab, anchor=east] at (33,24) {\iflangja{ポインタの表}{table of pointers}};
  \foreach \i/\x in {1/38, 2/46, 3/54} {
    \node[cell] (t\i) at (\x,24) {};
    \node[lab, anchor=south] at (\x,26.5) {\i};
  }
  \node[cell] at (62,24) {\ldots};
  \node[lab, anchor=west] at (67,24) {\iflangja{（添字 = スレッド ID）}{(index = thread ID)}};
  \fill (t1) circle (0.6);
  \fill (t2) circle (0.6);
  \draw[->] (t1.center) -- (54.5,9);
  \draw[->] (t2.center) .. controls (46,15) and (100,16) .. (106.5,9);
\end{tikzpicture}

  \caption{連続して割り当てられたユーザレベルスレッドのデータ構造．
    スレッド識別子は，構造へのポインタの表への添字である．T2 のスタックが
    あふれても，T1 の構造に達する前にレッドゾーンに入る．}
  \label{fig:l05-ult-layout}
\end{figure}

\needspace{13\baselineskip}
\begin{example}
  あるライブラリが各スレッドに \SI{64}{\kibi\byte} のブロックを与える．
  下端に \SI{4}{\kibi\byte} のレッドゾーン，\SI{56}{\kibi\byte} のスタック，
  上端にスレッド構造とそのスレッドローカルストレージのための
  \SI{4}{\kibi\byte} である．領域の先頭からの \si{\kibi\byte} 単位で測ると，
  ブロック $k$ は $[64k, 64k+64)$ を占める．したがってスレッド~3 は
  レッドゾーンを $[192, 196)$ に，スタックを $[196, 252)$ に，構造を
  $[252, 256)$ に持ち，スレッド~2 の構造は $[188, 192)$ にある．
  スレッド~3 が \SI{56}{\kibi\byte} を超えるスタックを必要とすると，次の
  プッシュは $[192, 196)$ に入ってフォールトする．レッドゾーンがなければ，
  スレッド~2 の構造を上書きしてしまう．1024 スレッドでは領域は
  \SI{64}{\mebi\byte} の仮想アドレス空間に及ぶが，そのうち
  \SI{4}{\mebi\byte} のレッドゾーンは物理メモリを使わない．\margin{x86 の
    NPTL も同じ配置を保つ．スレッド記述子はスタックブロックの上端にあり，
    glibc はスタックの下に既定で1ページ（\SI{4}{\kibi\byte}）のガード
    ページを置く．\texttt{pthread\_attr\_setguardsize} で変更できる．}
\end{example}

\subsection{カーネルレベルのデータ構造}

SunOS~5.0 のカーネルは，プロセスとスレッドについて4種類の構造を保持する
（\cref{tab:l05-kernel-structs}）．LWP はユーザレベルスレッドのデータ構造
に似ているが，カーネルから見える．その内容はプロセスが実行中の間だけ必要
なので，LWP はプロセスとともにスワップアウトできる．カーネルレベルスレッド
構造は，スケジューリング情報のように，プロセスが実行中でないときにも
カーネルが必要とする情報を保持し，決してスワップアウトされない．CPU 構造は，
各 CPU で現在実行中のスレッドを示す．\margin{予備のレジスタを持つ SPARC
プロセッサでは，SunOS は現在のスレッドへのポインタを専用のレジスタに保持
する．}

\begin{table}[htbp]
  \centering
  \caption{SunOS~5.0 におけるプロセスとスレッドのためのカーネルデータ構造．}
  \label{tab:l05-kernel-structs}
  \small
  \begin{tabular}{@{}>{\raggedright\arraybackslash}p{0.22\linewidth}>{\raggedright\arraybackslash}p{0.71\linewidth}@{}}
    \toprule
    構造 & 内容 \\
    \midrule
    プロセス & カーネルレベルスレッドのリスト，仮想アドレス空間，ユーザの
      資格情報，シグナルハンドラ \\
    LWP & ユーザレベルのレジスタ，システムコールの引数，資源使用情報，
      シグナルマスク．スワップ可能 \\
    カーネルレベルスレッド & カーネルレベルのレジスタ，スタックポインタ，
      スケジューリング情報，結び付けられた LWP・プロセス・CPU への
      ポインタ．スワップ不可 \\
    CPU & 現在のスレッド，カーネルレベルスレッドのリスト，ディスパッチと
      割り込み処理の情報 \\
    \bottomrule
  \end{tabular}
\end{table}

\section{ライブラリとカーネルの相互作用}
\label{sec:l05-interaction}

ユーザレベルライブラリはカーネルの中で何が起きているかを知らず，カーネルは
ライブラリの中で何が起きているかを知らない．\source{P2L4，スレッド管理の
基本的な相互作用，スレッド管理の可視性と設計；Stein と Shah
\cite{stein1992}；Silberschatz ら ch.~4．}システムコールと特別なシグナル
が，2つのレベルの協調を可能にする．アプリケーションは，自分のスレッドの
うち何個が同時に実行されるべきか---\texttt{pthread\_setconcurrency} で
設定する\emph{並行度}---をライブラリに伝え，ライブラリはシステムコールで
カーネルにカーネルレベルスレッドを要求し，足りなくなればさらに要求する．
カーネルは，カーネルレベルスレッドがブロックしたことなど，ライブラリの
スレッドに影響する事象をライブラリに通知できる．

\needspace{6\baselineskip}
\begin{example}
  あるプロセスが6つのユーザレベルスレッドを持ち，そのうち同時に実行される
  必要があるのは高々3つであるとする．\cref{tab:l05-interaction} は，
  ライブラリとカーネルがカーネルレベルスレッドの数をどう調整するかを
  たどったものである．\margin{\texttt{pthread\_setconcurrency} は Linux では
    何の効果も持たない．一対一のライブラリでは選ぶ余地がなく，POSIX も
    実装がこのヒントを無視することを許している．}
\end{example}

\begin{table}[htbp]
  \centering
  \caption{6つのユーザレベルスレッド（U1--U6）を持つプロセスにおける，
    ユーザレベルライブラリとカーネルの協調．}
  \label{tab:l05-interaction}
  \small
  \begin{tabular}{@{}c>{\raggedright\arraybackslash}p{0.63\linewidth}>{\raggedright\arraybackslash}p{0.18\linewidth}@{}}
    \toprule
    ステップ & 事象 & KLT \\
    \midrule
    1 & ライブラリがカーネルに3つのカーネルレベルスレッドを要求し，
        カーネルが K1--K3 を生成する． & K1--K3 \\
    2 & U1--U3 が K1--K3 上で実行され，U4--U6 はユーザレベルの条件で待つ．
      & 3 実行中 \\
    3 & U1 と U2 がブロッキング I/O 要求を発行し，K1 と K2 がカーネル内で
        ブロックする．その間に U4 が実行可能になる． &
        2 ブロック，1 実行中 \\
    4 & U3 もブロッキング要求を発行するので，U4 を実行できるカーネル
        レベルスレッドがなくなり，プロセス全体が停止してしまう．そこで
        カーネルは，K3 がブロックするときにライブラリに通知する
        ---SunOS~5.0 ではシグナル \texttt{SIGWAITING} を用いる．
        ライブラリはもう1つカーネルレベルスレッドを要求し，カーネルが
        K4 を生成して，その上で U4 が実行される． &
        3 ブロック，1 実行中 \\
    5 & I/O が完了し，U1--U3 が実行を続ける．その後 K4 は長い間アイドルの
        ままとなり，ライブラリはそれを破棄する． & K1--K3 \\
    \bottomrule
  \end{tabular}
\end{table}

\subsection{スレッド管理の可視性}

どちらのレベルも全体像の一部しか見ていない．
\begin{description}
  \item[カーネル] カーネルレベルスレッド，CPU，カーネルレベルスケジューラ
    の判断が見える．
  \item[ユーザレベルライブラリ] ユーザレベルスレッドと，自プロセスに
    割り当てられたカーネルレベルスレッドが見える．
\end{description}
ライブラリは，第3章のバウンドスレッドのように，ユーザレベルスレッドを
カーネルレベルスレッドに結び付けることができ，カーネルはカーネルレベル
スレッドを特定の CPU に固定することができる．どちらの判断も，もう一方の
レベルからは見えない．\margin{POSIX ではバウンドスレッドは
\texttt{PTHREAD\_SCOPE\_SYSTEM} であり，NPTL はこのスコープにしか対応
しない．カーネルレベルスレッドを CPU に固定するのは第7章のキャッシュ
アフィニティである．}

可視性の欠如は問題を引き起こす．あるユーザレベルスレッドがライブラリの
実装するミューテックスを獲得し，その後カーネルが，例えばタイムスライスが
尽きたために，そのスレッドが動いているカーネルレベルスレッドをプリエンプト
したとする．そのミューテックスを必要とする他のユーザレベルスレッドは，
たとえカーネルが，それらがスケジュールされているカーネルレベルスレッドを
実行していても，先に進めない．プリエンプトされたカーネルレベルスレッドを
カーネルが再び実行するまで何も変わらないが，ミューテックスが見えていない
カーネルには，それをすぐに実行する理由がないからである．多対多モデルでは，
ユーザレベルのスケジューリングの判断や，ユーザレベルスレッドのカーネルレベルスレッドへの
対応付けの変更は，ライブラリにしか見えない．一対一モデルはこれらの問題の
多くを避ける．カーネルがすべてのスレッドを見るからである．\sidenote{多対多
モデルに対する一般的な対策は，実行中のスレッド数を変える事象をカーネルが
すべてライブラリに知らせることである．Anderson ら \cite{anderson1992} は
このアップコールを\emph{スケジューラアクティベーション}と呼んだ．スレッドが
ブロックし，ブロックを解かれ，あるいはプリエンプトされるたびに，カーネルは
新しいアクティベーションをライブラリに渡し，ライブラリは直ちに再スケジュール
できる．\texttt{SIGWAITING} は同じ考え方の限定された形である．NetBSD と
FreeBSD はアクティベーションを実装したが，のちに削除した．多対多モデル自体
と同様に，得られるものより高くついたのである．}

ライブラリのスケジューラはプロセスのアドレス空間にある通常のコードである
から，実行がライブラリに入ったときにしか動かない．
\begin{itemize}
  \item ユーザレベルスレッドが明示的に実行権を譲る．
  \item ライブラリが設定したタイマが満了する．
  \item ユーザレベルスレッドが，ミューテックスのロックやアンロックなど，
    どのスレッドが実行可能かを変えうるライブラリ関数を呼ぶ．
  \item ブロックしていたユーザレベルスレッドが実行可能になる．
  \item 例えばタイマのために，カーネルからシグナルが届く．
\end{itemize}

\section{実装上の問題}
\label{sec:l05-issues}

マルチプロセッサはライブラリにさらなる問題を突きつけ，スレッド生成のコスト
はスレッドの再利用を促す．\source{P2L4，複数 CPU における問題，同期に
関する問題，スレッドの破棄；Eykholt ら \cite{eykholt1992}；OSTEP
ch.~28．}

\subsection{複数 CPU 上のユーザレベルスレッド}

マルチプロセッサでは，ライブラリは複数のカーネルレベルスレッド上で同時に
実行され，ある CPU で下した判断が別の CPU での行動を必要とすることがある．
優先度が $\text{T3} > \text{T2} > \text{T1}$ である3つのユーザレベル
スレッド T1，T2，T3 を，ライブラリが2つの CPU 上の2つのカーネルレベル
スレッド K1 と K2 で実行しているとする（\cref{fig:l05-multi-cpu}）．T2 は
CPU~1 でミューテックス \texttt{m} を保持している．T3 は CPU~2 で実行され，
\texttt{m} をロックしようとしてブロックし，ライブラリは K2 に T1 を
スケジュールする．T2 が \texttt{m} をアンロックすると，CPU~1 で実行されて
いるライブラリは，T3 が再び実行可能になり，実行中の T1 より優先度が高い
ことを知る．T1 をプリエンプトしなければならないが，CPU~1 上のライブラリの
コードは CPU~2 のレジスタを変更できない．そこでライブラリは K2 にシグナルを
送る．このシグナルによって K2 は CPU~2 でライブラリを実行し，ライブラリが
T1 をプリエンプトして T3 をスケジュールする．

\begin{figure}[htbp]
  \centering
  % Priority scheduling of user-level threads on two CPUs: when T2 unlocks the
% mutex, the library must signal the other kernel-level thread to preempt T1.
% Usage: % Priority scheduling of user-level threads on two CPUs: when T2 unlocks the
% mutex, the library must signal the other kernel-level thread to preempt T1.
% Usage: % Priority scheduling of user-level threads on two CPUs: when T2 unlocks the
% mutex, the library must signal the other kernel-level thread to preempt T1.
% Usage: \input{figures/l05-multi-cpu}   (width ~116 mm)
\begin{tikzpicture}[x=1mm, y=1mm, font=\footnotesize\sffamily,
  lab/.style={font=\scriptsize\sffamily, align=center},
  row/.style={font=\scriptsize\sffamily, anchor=east, align=right},
  >={Stealth[length=1.6mm]}]
  % \lfseg{x1}{x2}{y}{fill}{text}
  \def\lfseg#1#2#3#4#5{%
    \draw[fill=#4] (#1,#3-2.5) rectangle (#2,#3+2.5);
    \node[lab] at ({(#1+#2)/2},#3) {#5};}
  % event lines
  \draw[densely dashed, ln-muted] (44,24) -- (44,-2);
  \draw[densely dashed, ln-muted] (74,24) -- (74,-2);
  \node[lab, anchor=south, fill=white, inner sep=1pt] at (44,24)
    {\iflangja{T3 が m でブロック}{T3 blocks on m}};
  \node[lab, anchor=south, fill=white, inner sep=1pt] at (74,24)
    {\iflangja{T2 が m を解放}{T2 unlocks m}};
  % CPU 1
  \node[row] at (24,15) {CPU 1\\(KLT K1)};
  \lfseg{26}{74}{15}{ln-box}{\iflangja{T2（m を保持）}{T2 (holds m)}}
  \lfseg{74}{116}{15}{white}{T2}
  % CPU 2
  \node[row] at (24,4) {CPU 2\\(KLT K2)};
  \lfseg{26}{44}{4}{white}{T3}
  \lfseg{44}{48}{4}{black!35}{}
  \lfseg{48}{79}{4}{white}{T1}
  \lfseg{79}{83}{4}{black!35}{}
  \lfseg{83}{116}{4}{ln-box}{\iflangja{T3（m を保持）}{T3 (holds m)}}
  % signal from K1 to K2
  \draw[->, thick, ln-caution] (74,12.5) -- node[lab, right, pos=0.45, text=ln-caution]
    {\iflangja{シグナル}{signal}} (79,6.5);
  % legend
  \draw[fill=black!35] (26,-9) rectangle (30,-6);
  \node[lab, anchor=west] at (31,-7.5) {\iflangja{ライブラリのスケジューラ}{library scheduler}};
  \draw[fill=ln-box] (64,-9) rectangle (68,-6);
  \node[lab, anchor=west] at (69,-7.5) {\iflangja{m を保持}{m held}};
  \draw[->] (92,-7.5) -- node[lab, above] {\iflangja{時間}{time}} (116,-7.5);
\end{tikzpicture}
   (width ~116 mm)
\begin{tikzpicture}[x=1mm, y=1mm, font=\footnotesize\sffamily,
  lab/.style={font=\scriptsize\sffamily, align=center},
  row/.style={font=\scriptsize\sffamily, anchor=east, align=right},
  >={Stealth[length=1.6mm]}]
  % \lfseg{x1}{x2}{y}{fill}{text}
  \def\lfseg#1#2#3#4#5{%
    \draw[fill=#4] (#1,#3-2.5) rectangle (#2,#3+2.5);
    \node[lab] at ({(#1+#2)/2},#3) {#5};}
  % event lines
  \draw[densely dashed, ln-muted] (44,24) -- (44,-2);
  \draw[densely dashed, ln-muted] (74,24) -- (74,-2);
  \node[lab, anchor=south, fill=white, inner sep=1pt] at (44,24)
    {\iflangja{T3 が m でブロック}{T3 blocks on m}};
  \node[lab, anchor=south, fill=white, inner sep=1pt] at (74,24)
    {\iflangja{T2 が m を解放}{T2 unlocks m}};
  % CPU 1
  \node[row] at (24,15) {CPU 1\\(KLT K1)};
  \lfseg{26}{74}{15}{ln-box}{\iflangja{T2（m を保持）}{T2 (holds m)}}
  \lfseg{74}{116}{15}{white}{T2}
  % CPU 2
  \node[row] at (24,4) {CPU 2\\(KLT K2)};
  \lfseg{26}{44}{4}{white}{T3}
  \lfseg{44}{48}{4}{black!35}{}
  \lfseg{48}{79}{4}{white}{T1}
  \lfseg{79}{83}{4}{black!35}{}
  \lfseg{83}{116}{4}{ln-box}{\iflangja{T3（m を保持）}{T3 (holds m)}}
  % signal from K1 to K2
  \draw[->, thick, ln-caution] (74,12.5) -- node[lab, right, pos=0.45, text=ln-caution]
    {\iflangja{シグナル}{signal}} (79,6.5);
  % legend
  \draw[fill=black!35] (26,-9) rectangle (30,-6);
  \node[lab, anchor=west] at (31,-7.5) {\iflangja{ライブラリのスケジューラ}{library scheduler}};
  \draw[fill=ln-box] (64,-9) rectangle (68,-6);
  \node[lab, anchor=west] at (69,-7.5) {\iflangja{m を保持}{m held}};
  \draw[->] (92,-7.5) -- node[lab, above] {\iflangja{時間}{time}} (116,-7.5);
\end{tikzpicture}
   (width ~116 mm)
\begin{tikzpicture}[x=1mm, y=1mm, font=\footnotesize\sffamily,
  lab/.style={font=\scriptsize\sffamily, align=center},
  row/.style={font=\scriptsize\sffamily, anchor=east, align=right},
  >={Stealth[length=1.6mm]}]
  % \lfseg{x1}{x2}{y}{fill}{text}
  \def\lfseg#1#2#3#4#5{%
    \draw[fill=#4] (#1,#3-2.5) rectangle (#2,#3+2.5);
    \node[lab] at ({(#1+#2)/2},#3) {#5};}
  % event lines
  \draw[densely dashed, ln-muted] (44,24) -- (44,-2);
  \draw[densely dashed, ln-muted] (74,24) -- (74,-2);
  \node[lab, anchor=south, fill=white, inner sep=1pt] at (44,24)
    {\iflangja{T3 が m でブロック}{T3 blocks on m}};
  \node[lab, anchor=south, fill=white, inner sep=1pt] at (74,24)
    {\iflangja{T2 が m を解放}{T2 unlocks m}};
  % CPU 1
  \node[row] at (24,15) {CPU 1\\(KLT K1)};
  \lfseg{26}{74}{15}{ln-box}{\iflangja{T2（m を保持）}{T2 (holds m)}}
  \lfseg{74}{116}{15}{white}{T2}
  % CPU 2
  \node[row] at (24,4) {CPU 2\\(KLT K2)};
  \lfseg{26}{44}{4}{white}{T3}
  \lfseg{44}{48}{4}{black!35}{}
  \lfseg{48}{79}{4}{white}{T1}
  \lfseg{79}{83}{4}{black!35}{}
  \lfseg{83}{116}{4}{ln-box}{\iflangja{T3（m を保持）}{T3 (holds m)}}
  % signal from K1 to K2
  \draw[->, thick, ln-caution] (74,12.5) -- node[lab, right, pos=0.45, text=ln-caution]
    {\iflangja{シグナル}{signal}} (79,6.5);
  % legend
  \draw[fill=black!35] (26,-9) rectangle (30,-6);
  \node[lab, anchor=west] at (31,-7.5) {\iflangja{ライブラリのスケジューラ}{library scheduler}};
  \draw[fill=ln-box] (64,-9) rectangle (68,-6);
  \node[lab, anchor=west] at (69,-7.5) {\iflangja{m を保持}{m held}};
  \draw[->] (92,-7.5) -- node[lab, above] {\iflangja{時間}{time}} (116,-7.5);
\end{tikzpicture}

  \caption{2つの CPU 上の，優先度が $\text{T3} > \text{T2} > \text{T1}$ で
    あるユーザレベルスレッド．T2 が \texttt{m} をアンロックすると，CPU~1 の
    ライブラリは K2 にシグナルを送り，CPU~2 のライブラリが T1 をプリエンプト
    して T3 を実行する．}
  \label{fig:l05-multi-cpu}
\end{figure}

\subsection{適応型ミューテックス}

マルチプロセッサは，ミューテックスを待つ最良の方法も変える．CPU~1 である
スレッドがミューテックスを保持しているときに，CPU~2 で実行されている別の
スレッドがそれをロックしようとしたとする．通常，2番目のスレッドは
ミューテックスの待ちキューに置かれてブロックし，そのために今1回，
起こされるときにもう1回のコンテキストスイッチが必要になる．クリティカル
セクションが短く，その所有者が別の CPU で実行中であれば，これらの
コンテキストスイッチが完了するより前に所有者がミューテックスを解放する
可能性が高く，2番目のスレッドは\emph{スピン}する方が---自分の CPU に
とどまり，ミューテックスが空いたかどうかを繰り返し調べる方が---安上がり
である．所有者が実行中でなければ，所有者が再びスケジュールされるまで
ミューテックスは解放されえないので，スピンは無意味である．

\begin{definition}[適応型ミューテックス]
  \term[てきおうがたみゅーてっくす]{適応型ミューテックス}[adaptive mutex]%
  とは，そのロック操作が，ミューテックスの所有者が別の CPU で実行中の間は
  スピンし，そうでなければ呼び出したスレッドをブロックするミューテックス
  である．
\end{definition}

適応型ミューテックスが意味を持つのはマルチプロセッサに限られ---単一の CPU
では，別のスレッドがスピンしている間に所有者が実行中であることはありえない
---かつ短いクリティカルセクションに限られる．長いクリティカルセクションに
はブロックの方が適する．スピンとブロックのどちらを選ぶかを決めるには，
ロック操作が所有者の実行中かどうかを知らなければならないので，
ミューテックスのデータ構造はその所有者を記録する．スピンロックの基礎と
してのスピンは第10章で改めて扱う．\margin{OSTEP では，しばらくスピンして
から眠るロックを\emph{2相ロック}と呼んでいる．Linux カーネルのミューテックス
も，ここと同じく所有者が実行中の間だけスピンする．}

\subsection{スレッドの破棄}

スレッドの生成は，そのデータ構造とスタックを割り当てて初期化することであり，
破棄はそれらを解放することである．アプリケーションはスレッドを頻繁に生成・
破棄するので，これらの構造を再利用することには利点がある．スレッドが終了
しても，その構造はすぐには解放されず，スレッドは\emph{デスロー}（death
row）と呼ばれるリストに置かれる．\term[りーぱすれっど]{リーパスレッド}%
[reaper thread]が定期的に実行され，デスロー上のスレッドを破棄してその
データ構造を解放する．リーパが実行される前にスレッド生成の要求が届いた
場合には，デスロー上のスレッドの構造とスタックが再利用され，それらを
割り当てて初期化するコストが節約される．\margin{glibc も NPTL のスレッドに
対して同じことを行う．解放されたスタックは最大 \SI{40}{\mebi\byte} の
キャッシュ（チューナブル \texttt{glibc.pthread.stack\_cache\_size}）に残り，
次の \texttt{pthread\_create} に渡される．}

\begin{keypoint}
  カーネルとユーザレベルライブラリは，それぞれ自分のスレッドしか見ない．
  システムコールとシグナルによって両者は協調できるが，ロックの所有者を
  プリエンプトする，あるいは別の CPU 上のスレッドをプリエンプトすると
  いった判断には，レベルをまたぐ明示的な通知が必要である．複数の CPU は
  スレッドの待ち方も変え---適応型ミューテックスは所有者が実行中の間だけ
  スピンする---終了したスレッドは直ちに解放されるのではなく再利用のために
  保持される．
\end{keypoint}

\section{割り込みとシグナル}
\label{sec:l05-int-sig}

スレッドを通常の実行から逸らす事象には2種類ある．割り込みとシグナルで
ある．\source{P2L4，割り込みとシグナルの導入；OSTEP ch.~6；Kerrisk
\cite{kerrisk2010} ch.~20．}シグナルは第1章で，オペレーティングシステムが
アプリケーションに渡す通知として導入した．割り込みは，ハードウェアが CPU に
届ける，これに対応する通知である．

\begin{definition}[割り込み]
  \term[わりこみ]{割り込み}[interrupt]とは，I/O デバイス，タイマ，他の CPU
  といった構成要素によって CPU の外部で生成され，何らかの条件が対応を要する
  ことを CPU に通知する事象である．
\end{definition}

どの割り込みが起こりうるかは物理的なプラットフォームで決まり，割り込みは
CPU が何を実行しているかとは無関係に非同期に現れる．一方シグナルは，CPU と
その上で動くソフトウェアが引き起こす事象である．どのシグナルが存在するかは
オペレーティングシステムが決め，シグナルは同期的にも非同期的にも現れうる．
とはいえ，2つの機構には共通点が多い（\cref{tab:l05-int-vs-sig}）．どちらも，
ハードウェアまたはオペレーティングシステムが定めた一意な識別子を持ち，
どちらもマスクできる，すなわち無効化・抑止できる．そして有効であれば，
どちらもハンドラを起動する．割り込みハンドラはオペレーティングシステムが
システム全体に対して設定し，シグナルハンドラは各プロセスが自分のために
設定する．

\begin{table}[htbp]
  \centering
  \caption{割り込みとシグナル．}
  \label{tab:l05-int-vs-sig}
  \small
  \begin{tabular}{@{}>{\raggedright\arraybackslash}p{0.22\linewidth}>{\raggedright\arraybackslash}p{0.345\linewidth}>{\raggedright\arraybackslash}p{0.345\linewidth}@{}}
    \toprule
    & 割り込み & シグナル \\
    \midrule
    生成元 & デバイス，タイマ，他の CPU & CPU とその上で動くソフトウェア \\
    定める主体 & ハードウェアプラットフォーム & オペレーティングシステム \\
    発生 & 非同期 & 同期または非同期 \\
    マスク & CPU ごと & 実行コンテキストごと \\
    ハンドラの設定 & オペレーティングシステムがシステム全体に対して
      & 各プロセスが自分のために \\
    \bottomrule
  \end{tabular}
\end{table}

\section{割り込みとシグナルの処理}
\label{sec:l05-handling}

デバイスは対応を必要とするとき，CPU との間をつなぐ PCI Express などの
インターコネクトを通じて CPU に割り込みを送る．\source{P2L4，割り込みの
処理，シグナルの処理；OSTEP ch.~6，36；Kerrisk ch.~20--21．}初期の
システムは割り込みごとに専用の物理線を用いていたが，現代のデバイスは
インターコネクト自体を通じてメッセージとして割り込みを送り，メッセージの
内容が割り込みを，したがってデバイスを識別する．専用の線を駆動する代わりに
このようなメッセージを送って発生させる割り込みを，
\term[めっせーじしぐなるわりこみ]{メッセージシグナル割り込み}[message signaled interrupt]%
\index{MSI|see{メッセージシグナル割り込み}}（MSI）と呼ぶ．

\begin{definition}[割り込みハンドラ]
  \term[わりこみはんどら]{割り込みハンドラ}[interrupt handler]とは，特定の
  割り込みに応じてオペレーティングシステムが実行するルーチンである．
  \term[わりこみはんどらてーぶる]{割り込みハンドラテーブル}[interrupt handler table]%
  は，各割り込み番号をそのハンドラの先頭アドレスに対応付ける．
\end{definition}

CPU は割り込み番号で割り込みハンドラテーブルを引き，そこで見つけた先頭
アドレスをプログラムカウンタに設定してハンドラを実行する
（\cref{fig:l05-int-sig}a）．どの割り込みが起こるかはハードウェアに固有で
あり，それらをどう処理するかはオペレーティングシステムに固有である．
\margin{この表は割り込みベクタテーブルとも呼ばれ，x86 では割り込み
ディスクリプタテーブルである．}

\begin{figure}[htbp]
  \centering
  % Handling an interrupt (system-wide handler table) and a signal (per-process
% handler table).
% Usage: % Handling an interrupt (system-wide handler table) and a signal (per-process
% handler table).
% Usage: % Handling an interrupt (system-wide handler table) and a signal (per-process
% handler table).
% Usage: \input{figures/l05-int-sig}   (width ~118 mm)
\begin{tikzpicture}[x=1mm, y=1mm, font=\scriptsize\sffamily,
  >={Stealth[length=1.5mm]},
  box/.style={draw, fill=white, minimum height=9mm, inner sep=1pt, align=center},
  act/.style={draw, fill=ln-accent!22, minimum height=9mm, inner sep=1pt,
              align=center},
  lab/.style={font=\scriptsize\sffamily, align=center},
  ttl/.style={font=\footnotesize\sffamily\bfseries, anchor=west},
  tcap/.style={font=\scriptsize\sffamily, text=ln-muted, anchor=south}]
  % \lstable{y-top}{row1a}{row1b}{row3a}{row3b}: three-row handler table at x=52..88
  \def\lstable#1#2#3#4#5{%
    \draw[fill=ln-box] (52,#1) rectangle (88,#1-15);
    \draw (52,#1-5) -- (88,#1-5);
    \draw (52,#1-10) -- (88,#1-10);
    \draw (65,#1) -- (65,#1-15);
    \node at (58.5,#1-2.5) {#2};  \node at (76.5,#1-2.5) {#3};
    \node at (58.5,#1-7.5) {\ldots}; \node at (76.5,#1-7.5) {\ldots};
    \node at (58.5,#1-12.5) {#4}; \node at (76.5,#1-12.5) {#5};}
  % ---- (a) interrupt
  \node[ttl] at (0,33) {\iflangja{(a) 割り込み}{(a) interrupt}};
  \node[box, minimum width=16mm] (dev) at (8,20) {\iflangja{デバイス}{device}};
  \node[act, minimum width=12mm] (cpu) at (38,20) {CPU};
  \draw[->, thick] (dev.east) -- node[lab, above] {MSI} node[lab, below]
    {\iflangja{（PCIe）}{(PCIe)}} (cpu.west);
  \lstable{27.5}{INT-1}{\iflangja{ハンドラ 1 の番地}{addr.\ handler 1}}%
    {INT-N}{\iflangja{ハンドラ N の番地}{addr.\ handler N}}
  \node[tcap] at (70,27.8) {\iflangja{割り込みハンドラテーブル（システム全体）}{interrupt handler table (system-wide)}};
  \draw[->] (cpu.east) -- node[lab, above] {N} (47,20) |- (52,15);
  \node[act, minimum width=20mm] (h1) at (107,15) {\iflangja{ハンドラ N}{handler N}};
  \draw[->, thick] (88,15) -- node[lab, above] {\iflangja{PC を}{set}}
    node[lab, below] {\iflangja{設定}{PC}} (h1.west);
  % ---- (b) signal
  \node[ttl] at (0,1) {\iflangja{(b) シグナル}{(b) signal}};
  \node[box, minimum width=16mm] (thr) at (8,-12) {\iflangja{P の\\スレッド}{thread\\of P}};
  \node[act, minimum width=12mm] (ker) at (38,-12) {\iflangja{カーネル}{kernel}};
  \draw[->, thick] (thr.east) -- node[lab, above] {\iflangja{不正な}{illegal}}
    node[lab, below] {\iflangja{アクセス}{access}} (ker.west);
  \lstable{-4.5}{SIG-1}{\iflangja{ハンドラ / 既定}{handler / default}}%
    {SIGSEGV}{\iflangja{ハンドラの番地}{addr.\ handler}}
  \node[tcap] at (70,-4.2) {\iflangja{P のシグナルハンドラテーブル}{signal handler table of P}};
  \draw[->] (ker.east) -- node[lab, above] {11} (47,-12) |- (52,-17);
  \node[act, minimum width=20mm] (h2) at (107,-17) {\iflangja{ハンドラ\\（スレッドの\\文脈で実行）}{handler\\(in thread's\\context)}};
  \draw[->, thick] (88,-17) -- node[lab, above] {\iflangja{PC を}{set}}
    node[lab, below] {\iflangja{設定}{PC}} (h2.west);
\end{tikzpicture}
   (width ~118 mm)
\begin{tikzpicture}[x=1mm, y=1mm, font=\scriptsize\sffamily,
  >={Stealth[length=1.5mm]},
  box/.style={draw, fill=white, minimum height=9mm, inner sep=1pt, align=center},
  act/.style={draw, fill=ln-accent!22, minimum height=9mm, inner sep=1pt,
              align=center},
  lab/.style={font=\scriptsize\sffamily, align=center},
  ttl/.style={font=\footnotesize\sffamily\bfseries, anchor=west},
  tcap/.style={font=\scriptsize\sffamily, text=ln-muted, anchor=south}]
  % \lstable{y-top}{row1a}{row1b}{row3a}{row3b}: three-row handler table at x=52..88
  \def\lstable#1#2#3#4#5{%
    \draw[fill=ln-box] (52,#1) rectangle (88,#1-15);
    \draw (52,#1-5) -- (88,#1-5);
    \draw (52,#1-10) -- (88,#1-10);
    \draw (65,#1) -- (65,#1-15);
    \node at (58.5,#1-2.5) {#2};  \node at (76.5,#1-2.5) {#3};
    \node at (58.5,#1-7.5) {\ldots}; \node at (76.5,#1-7.5) {\ldots};
    \node at (58.5,#1-12.5) {#4}; \node at (76.5,#1-12.5) {#5};}
  % ---- (a) interrupt
  \node[ttl] at (0,33) {\iflangja{(a) 割り込み}{(a) interrupt}};
  \node[box, minimum width=16mm] (dev) at (8,20) {\iflangja{デバイス}{device}};
  \node[act, minimum width=12mm] (cpu) at (38,20) {CPU};
  \draw[->, thick] (dev.east) -- node[lab, above] {MSI} node[lab, below]
    {\iflangja{（PCIe）}{(PCIe)}} (cpu.west);
  \lstable{27.5}{INT-1}{\iflangja{ハンドラ 1 の番地}{addr.\ handler 1}}%
    {INT-N}{\iflangja{ハンドラ N の番地}{addr.\ handler N}}
  \node[tcap] at (70,27.8) {\iflangja{割り込みハンドラテーブル（システム全体）}{interrupt handler table (system-wide)}};
  \draw[->] (cpu.east) -- node[lab, above] {N} (47,20) |- (52,15);
  \node[act, minimum width=20mm] (h1) at (107,15) {\iflangja{ハンドラ N}{handler N}};
  \draw[->, thick] (88,15) -- node[lab, above] {\iflangja{PC を}{set}}
    node[lab, below] {\iflangja{設定}{PC}} (h1.west);
  % ---- (b) signal
  \node[ttl] at (0,1) {\iflangja{(b) シグナル}{(b) signal}};
  \node[box, minimum width=16mm] (thr) at (8,-12) {\iflangja{P の\\スレッド}{thread\\of P}};
  \node[act, minimum width=12mm] (ker) at (38,-12) {\iflangja{カーネル}{kernel}};
  \draw[->, thick] (thr.east) -- node[lab, above] {\iflangja{不正な}{illegal}}
    node[lab, below] {\iflangja{アクセス}{access}} (ker.west);
  \lstable{-4.5}{SIG-1}{\iflangja{ハンドラ / 既定}{handler / default}}%
    {SIGSEGV}{\iflangja{ハンドラの番地}{addr.\ handler}}
  \node[tcap] at (70,-4.2) {\iflangja{P のシグナルハンドラテーブル}{signal handler table of P}};
  \draw[->] (ker.east) -- node[lab, above] {11} (47,-12) |- (52,-17);
  \node[act, minimum width=20mm] (h2) at (107,-17) {\iflangja{ハンドラ\\（スレッドの\\文脈で実行）}{handler\\(in thread's\\context)}};
  \draw[->, thick] (88,-17) -- node[lab, above] {\iflangja{PC を}{set}}
    node[lab, below] {\iflangja{設定}{PC}} (h2.west);
\end{tikzpicture}
   (width ~118 mm)
\begin{tikzpicture}[x=1mm, y=1mm, font=\scriptsize\sffamily,
  >={Stealth[length=1.5mm]},
  box/.style={draw, fill=white, minimum height=9mm, inner sep=1pt, align=center},
  act/.style={draw, fill=ln-accent!22, minimum height=9mm, inner sep=1pt,
              align=center},
  lab/.style={font=\scriptsize\sffamily, align=center},
  ttl/.style={font=\footnotesize\sffamily\bfseries, anchor=west},
  tcap/.style={font=\scriptsize\sffamily, text=ln-muted, anchor=south}]
  % \lstable{y-top}{row1a}{row1b}{row3a}{row3b}: three-row handler table at x=52..88
  \def\lstable#1#2#3#4#5{%
    \draw[fill=ln-box] (52,#1) rectangle (88,#1-15);
    \draw (52,#1-5) -- (88,#1-5);
    \draw (52,#1-10) -- (88,#1-10);
    \draw (65,#1) -- (65,#1-15);
    \node at (58.5,#1-2.5) {#2};  \node at (76.5,#1-2.5) {#3};
    \node at (58.5,#1-7.5) {\ldots}; \node at (76.5,#1-7.5) {\ldots};
    \node at (58.5,#1-12.5) {#4}; \node at (76.5,#1-12.5) {#5};}
  % ---- (a) interrupt
  \node[ttl] at (0,33) {\iflangja{(a) 割り込み}{(a) interrupt}};
  \node[box, minimum width=16mm] (dev) at (8,20) {\iflangja{デバイス}{device}};
  \node[act, minimum width=12mm] (cpu) at (38,20) {CPU};
  \draw[->, thick] (dev.east) -- node[lab, above] {MSI} node[lab, below]
    {\iflangja{（PCIe）}{(PCIe)}} (cpu.west);
  \lstable{27.5}{INT-1}{\iflangja{ハンドラ 1 の番地}{addr.\ handler 1}}%
    {INT-N}{\iflangja{ハンドラ N の番地}{addr.\ handler N}}
  \node[tcap] at (70,27.8) {\iflangja{割り込みハンドラテーブル（システム全体）}{interrupt handler table (system-wide)}};
  \draw[->] (cpu.east) -- node[lab, above] {N} (47,20) |- (52,15);
  \node[act, minimum width=20mm] (h1) at (107,15) {\iflangja{ハンドラ N}{handler N}};
  \draw[->, thick] (88,15) -- node[lab, above] {\iflangja{PC を}{set}}
    node[lab, below] {\iflangja{設定}{PC}} (h1.west);
  % ---- (b) signal
  \node[ttl] at (0,1) {\iflangja{(b) シグナル}{(b) signal}};
  \node[box, minimum width=16mm] (thr) at (8,-12) {\iflangja{P の\\スレッド}{thread\\of P}};
  \node[act, minimum width=12mm] (ker) at (38,-12) {\iflangja{カーネル}{kernel}};
  \draw[->, thick] (thr.east) -- node[lab, above] {\iflangja{不正な}{illegal}}
    node[lab, below] {\iflangja{アクセス}{access}} (ker.west);
  \lstable{-4.5}{SIG-1}{\iflangja{ハンドラ / 既定}{handler / default}}%
    {SIGSEGV}{\iflangja{ハンドラの番地}{addr.\ handler}}
  \node[tcap] at (70,-4.2) {\iflangja{P のシグナルハンドラテーブル}{signal handler table of P}};
  \draw[->] (ker.east) -- node[lab, above] {11} (47,-12) |- (52,-17);
  \node[act, minimum width=20mm] (h2) at (107,-17) {\iflangja{ハンドラ\\（スレッドの\\文脈で実行）}{handler\\(in thread's\\context)}};
  \draw[->, thick] (88,-17) -- node[lab, above] {\iflangja{PC を}{set}}
    node[lab, below] {\iflangja{設定}{PC}} (h2.west);
\end{tikzpicture}

  \caption{(a) デバイスが発生させた割り込みは，システム全体の割り込み
    ハンドラテーブルからハンドラを選ぶ．(b) 不正なメモリアクセスにより
    カーネルが \texttt{SIGSEGV} を発生させ，これがプロセスのシグナル
    ハンドラテーブルからハンドラを選ぶ．}
  \label{fig:l05-int-sig}
\end{figure}

シグナルはオペレーティングシステムが生成する．例えばスレッドが自分の
プロセスに割り当てられていないメモリにアクセスすると，オペレーティング
システムは \texttt{SIGSEGV}（Linux ではシグナル~11）を発生させる．%
\margin{Linux は 31 個の標準シグナルと，ひとまとまりのリアルタイムシグナルを
定める．glibc は後者の最初の2つを NPTL のために予約するので，
\texttt{SIGRTMIN} は 34，\texttt{SIGRTMAX} は 64 である
（\texttt{signal(7)}）．}
オペレーティングシステムはプロセスごとに，各シグナル番号を取るべき動作に
対応付ける\emph{シグナルハンドラテーブル}を保持する
（\cref{fig:l05-int-sig}b）．シグナルの集合はオペレーティングシステムに
固有であり，POSIX は標準的な集合を定めている．すべてのシグナルには既定の
動作があり，それはプロセスの終了，シグナルの無視，プロセスの終了とコア
ダンプ，プロセスの停止，停止したプロセスの再開のいずれかである．

\begin{definition}[シグナルハンドラ]
  \term[しぐなるはんどら]{シグナルハンドラ}[signal handler]とは，特定の
  シグナルが配送されたときに，それを受け取るスレッドの文脈で実行される
  ようプロセスが設定する関数である．
\end{definition}

プロセスはハンドラを \texttt{signal} で，あるいはより移植性の高い
\texttt{sigaction} で設定する．ほとんどのシグナルはこの方法で捕捉できるが，
いくつかはできない．特に \texttt{SIGKILL} と \texttt{SIGSTOP} には，常に
既定の動作が適用される．

\begin{example}
  \texttt{signal.h} と \texttt{unistd.h} を用いて，タイマが満了したときに
  作業ループを止めるプログラム．\caution{ハンドラが代入してよいのは
    \texttt{volatile sig\_atomic\_t} のオブジェクトだけであり，呼んでよいのは
    \texttt{signal-safety(7)} に列挙された関数だけである．\texttt{printf} や
    \texttt{malloc} はそこに含まれない．}
  \begin{lstlisting}[language=C, numbers=none]
static volatile sig_atomic_t expired = 0;

static void on_alarm(int signo) {
    (void)signo;
    expired = 1;            /* ハンドラは最小限に保つ */
}

int main(void) {
    struct sigaction sa;
    sa.sa_handler = on_alarm;
    sigemptyset(&sa.sa_mask);
    sa.sa_flags = 0;
    sigaction(SIGALRM, &sa, NULL);
    alarm(2);               /* 2 秒後に SIGALRM */
    while (!expired)
        do_work();
    return 0;
}
\end{lstlisting}
\end{example}

シグナルは，それを受け取るスレッドとの関係によって分類される．

\begin{definition}[同期シグナルと非同期シグナル]
  \term[どうきしぐなる]{同期シグナル}[synchronous signal]とは，特定の
  スレッドの動作に応じて発生し，そのスレッドに配送されるシグナルである．
  保護されたメモリへのアクセスに対する \texttt{SIGSEGV} や，ゼロ除算に
  対する \texttt{SIGFPE} がその例である．
  \term[ひどうきしぐなる]{非同期シグナル}[asynchronous signal]とは，
  受け取るスレッドが何を実行しているかとは無関係に発生するシグナルである．
  他のプロセスが \texttt{kill} で送る \texttt{SIGKILL} や，タイマの満了に
  よる \texttt{SIGALRM} がその例である．
\end{definition}

シグナルは，例えば \texttt{pthread\_kill} で特定のスレッドに明示的に送る
こともできる．その場合，同期シグナルと同様に，そのスレッドにのみ配送される．

\section{割り込みとシグナルの無効化}
\label{sec:l05-masks}

割り込みハンドラとシグナルハンドラは，割り込まれたスレッドの文脈で実行
される．そのスレッドのプログラムカウンタがハンドラに設定され，ハンドラは
スケジューラを介さずにそのスレッドの代わりに実行される．このときシグナル
ハンドラはスレッド自身のスタックを，割り込みハンドラはカーネルスタックを
用いる．\source{P2L4，なぜ割り込みやシグナルを無効化するのか，シグナル
マスクの詳細，マルチコアシステムにおける割り込み，シグナルの種類；Kerrisk
ch.~20--22．}このことは，ハンドラが共有状態にアクセスするときに問題を生む．
割り込まれたスレッドがミューテックス \texttt{m} を保持していて，ハンドラが
\texttt{m} をロックしようとすると，ハンドラはブロックする．\texttt{m} の
所有者は割り込まれたスレッド自身であり，そのスレッドはハンドラが戻るまで
先に進めない．こうしてスレッドは自分自身のハンドラとデッドロックする．

これを避ける方法は2つある．1つは，ハンドラのコードを単純に保ち，決して
ミューテックスをロックしないようにすることであるが，これはしばしば制限が
強すぎる．もう1つは，ハンドラが必要とするミューテックスを保持している間は
該当する割り込みやシグナルを無効化することで，ハンドラがいつ実行されうるか
をコードの側で制御することである．

\begin{definition}[割り込みマスクとシグナルマスク]
  \term[わりこみますく]{割り込みマスク}[interrupt mask]および
  \term[しぐなるますく]{シグナルマスク}[signal mask]とは，割り込みまたは
  シグナル1つにつき1ビットを持つビット列であり，そのビットの値が，その
  割り込みまたはシグナルが有効（1）か無効（0）かを示す．
\end{definition}

\needspace{7\baselineskip}
クリティカルセクションの前後でシグナルを無効化するスレッドは，次の形を
とる．
\begin{lstlisting}[language=C, numbers=none]
clear_field_in_mask(mask);  // 無効化
lock(mutex) {
  ...                       // ここではハンドラは実行されない
} // unlock
reset_field_in_mask(mask);  // 有効化
\end{lstlisting}
無効化されている間に発生した割り込みやシグナルは失われない．それは
\emph{保留}のまま残り，マスクが再びそれを有効にすると直ちにそのハンドラが
実行される．無効化されている間に同じシグナルが複数回発生しても，通常ハンドラは
1回しか実行されない．シグナルを複数回発生させても，ハンドラが複数回実行
される保証はない（\cref{sec:l05-types}）．\caution{POSIX のシグナル集合は
\emph{ブロックされる}シグナルを列挙する．集合の要素であるシグナルは無効で
あり，上のビットの規約とは逆である．}

\subsection{マスクの設定}

割り込みマスクは CPU ごとに保持される．ある CPU のマスクが割り込みを無効に
していれば，ハードウェアの割り込み配送機構はその割り込みをその CPU に
届けない．カーネルはハードウェアに書き込むことでマスクを設定する．シグナル
マスクは実行コンテキストごとに保持される．カーネルは，マスクがそのシグナルを
無効にしているカーネルレベルスレッドには割り込まない．そして，ユーザレベル
スレッドがカーネルレベルスレッドの上で実行される場合，カーネルには
カーネルレベルのマスクしか見えない．これについては
\cref{sec:l05-thread-signals}で改めて扱う．シングルスレッドのプロセスは
\texttt{sigprocmask} で，マルチスレッドプロセスのスレッドは
\texttt{pthread\_sigmask} でシグナルマスクを変更する．どちらも
\texttt{SIG\_BLOCK}，\texttt{SIG\_UNBLOCK}，\texttt{SIG\_SETMASK} のいずれ
かの操作を取る．

\needspace{11\baselineskip}
\begin{example}
  上の形を POSIX の呼び出しで書くと，同じミューテックスをロックする
  \texttt{SIGALRM} のハンドラに対して次のようになる．
  \begin{lstlisting}[language=C, numbers=none]
sigset_t set, old;
sigemptyset(&set);
sigaddset(&set, SIGALRM);
pthread_sigmask(SIG_BLOCK, &set, &old);    // 無効化
pthread_mutex_lock(&m);
/* ... クリティカルセクション ... */
pthread_mutex_unlock(&m);
pthread_sigmask(SIG_SETMASK, &old, NULL);  // 復元
\end{lstlisting}
  \caution{\texttt{pthread\_mutex\_lock} は非同期シグナル安全ではない．
    ハンドラはフラグを立てるにとどめ，それを実行しうるすべてのスレッドが
    そのシグナルをブロックしなければならない．}
\end{example}

\subsection{マルチコアシステムにおける割り込み}

マルチコアシステムでは，割り込みはそれを有効にしているどの CPU にも
届けられうる．ある割り込みを1つのコアだけで有効にし，そのコアだけが処理
するようにすれば，他のすべてのコアでその処理が引き起こすオーバーヘッドと
乱れを避けられる．

\needspace{5\baselineskip}
\begin{keypoint}
  ハンドラは割り込まれたスレッドの文脈で実行される．ハンドラが必要とする
  ミューテックスを保持するコードは，それを保持している間，対応する割り込み
  またはシグナルを無効化しなければならない．その間に発生した事象は，再び
  有効になるまで保留のまま残る．
\end{keypoint}

\subsection{シグナルの種類}
\label{sec:l05-types}

シグナルは，ハンドラが実行される前に同じシグナルが複数回発生したときに何が
起こるかで異なる．
\begin{description}
  \item[ワンショットシグナル]
    \term[わんしょっとしぐなる]{ワンショットシグナル}[one-shot signal]では，
    保留中の $n$ 回分の発生は1回分と等価である．ハンドラは少なくとも1回は
    実行されることが保証されるが，$n$ 回実行されるとは限らない．システムに
    よっては，ハンドラは実行後に既定の動作に戻されるので，プロセスは再び
    そのシグナルを捕捉するために明示的に設定し直さなければならない．
  \item[リアルタイムシグナル]
    \term[りあるたいむしぐなる]{リアルタイムシグナル}[real-time signal]では，
    個々の発生がキューに入れられる．シグナルが $n$ 回発生すれば，ハンドラは
    $n$ 回呼ばれる．POSIX はこれらのために \texttt{SIGRTMIN} から
    \texttt{SIGRTMAX} までの範囲を予約している．
\end{description}

\begin{example}
  あるスレッドが，標準のシグナルである \texttt{SIGUSR1} と，リアルタイム
  シグナルである \texttt{SIGRTMIN} の両方をブロックする．ブロックされて
  いる間に各シグナルが4回発生し，その後スレッドはそれらのブロックを解除
  する．\texttt{SIGUSR1} のハンドラは，4回分の発生が1個の保留シグナルに
  まとまるので1回実行され，\texttt{SIGRTMIN} のハンドラは4回実行される．
\end{example}

\section{スレッドとしての割り込み}
\label{sec:l05-int-threads}

SunOS~5.0 は，スレッドとそのハンドラの間のデッドロックを別の方法で避けて
いる．割り込みをスレッドとして処理するのである．\source{P2L4，スレッドと
しての割り込み，割り込みをスレッドとする方式の性能；Eykholt ら
\cite{eykholt1992}．}割り込みハンドラが独自のスレッドとして実行されるなら，
それは他のスレッドと同様にミューテックスでブロックできる．所有者---それは
割り込まれたスレッドかもしれない---がミューテックスを解放するまで待ち，
その後スケジューラがそれを実行する．割り込みごとに新しいスレッドを生成するのは
コストが高すぎるので，割り込みスレッドはあらかじめ生成・初期化しておき，
そのうちの1つが完全なスレッドになるかどうかは実行時に決める．割り込みが
起こると，そのようにあらかじめ割り当てられたスレッドのスタックに切り替わり，
割り込まれたスレッドはその下にピン留めされたままとなって実行できない．
ブロックしないハンドラはそこで終わり，ピン留めされたスレッドが再開する．
ハンドラがミューテックスでブロックした場合に限り，割り込みスレッドは完全なスレッドと
なり，スケジューラはそれを他のスレッドと同様に扱って，ピン留めされていた
スレッドは直ちに実行を続ける．

同じ考え方は，割り込み処理の作業を分割する一般的な方法の基礎にもなっている
（\cref{fig:l05-top-bottom}）．\margin{Linux の遅延実行機構のうち，
ブロックしてよいのはワークキューとスレッド化ハンドラだけである．}

\begin{definition}[トップハーフとボトムハーフ]
  割り込み処理の\term[とっぷはーふ]{トップハーフ}[top half]は，割り込みが
  発生すると直ちに実行される部分であり，最小限の処理だけを行い，高速で
  ブロックしないものでなければならない．
  \term[ぼとむはーふ]{ボトムハーフ}[bottom half]は遅延される部分であり，
  任意の複雑さを持ちうる．スレッドがそれを実行する場合，それは他のスレッド
  と同様にスケジュールされ，ブロックしてもよい．
\end{definition}

\begin{figure}[htbp]
  \centering
  % Interrupt handling split into a top half, executed immediately, and a
% bottom half, executed later like a thread.
% Usage: % Interrupt handling split into a top half, executed immediately, and a
% bottom half, executed later like a thread.
% Usage: % Interrupt handling split into a top half, executed immediately, and a
% bottom half, executed later like a thread.
% Usage: \input{figures/l05-top-bottom}   (width ~116 mm)
\begin{tikzpicture}[x=1mm, y=1mm, font=\footnotesize\sffamily,
  lab/.style={font=\scriptsize\sffamily, align=center},
  row/.style={font=\scriptsize\sffamily, anchor=east},
  >={Stealth[length=1.6mm]}]
  \def\ltbseg#1#2#3#4#5{%
    \draw[fill=#4] (#1,#3-2.5) rectangle (#2,#3+2.5);
    \node[lab] at ({(#1+#2)/2},#3) {#5};}
  \node[row] at (14,10) {CPU};
  \ltbseg{16}{44}{10}{white}{\iflangja{スレッド T}{thread T}}
  \ltbseg{44}{61}{10}{ln-accent!22}{\iflangja{トップハーフ}{top half}}
  \ltbseg{61}{80}{10}{white}{\iflangja{スレッド T}{thread T}}
  \ltbseg{80}{104}{10}{ln-box}{\iflangja{ボトムハーフ}{bottom half}}
  \ltbseg{104}{116}{10}{white}{T}
  % interrupt arrival
  \draw[->, thick, ln-caution] (44,20) -- (44,12.5);
  \node[lab, anchor=south, text=ln-caution] at (44,20) {\iflangja{割り込み}{interrupt}};
  % top half defers the bottom half
  \draw[->, ln-accent] (52.5,7.5) .. controls (56,1) and (86,1) .. (92,7.5);
  \node[lab, text=ln-accent] at (72,4.8) {\iflangja{後で実行するよう登録}{deferred}};
  % annotations
  \node[lab, anchor=north] at (52.5,-1)
    {\iflangja{直ちに実行\\高速・ブロックしない}{runs immediately;\\fast, non-blocking}};
  \node[lab, anchor=north] at (92,-1)
    {\iflangja{スレッドとしてスケジュール\\ブロック可能}{scheduled like a thread;\\may block}};
  \draw[->] (92,20) -- node[lab, above] {\iflangja{時間}{time}} (116,20);
\end{tikzpicture}
   (width ~116 mm)
\begin{tikzpicture}[x=1mm, y=1mm, font=\footnotesize\sffamily,
  lab/.style={font=\scriptsize\sffamily, align=center},
  row/.style={font=\scriptsize\sffamily, anchor=east},
  >={Stealth[length=1.6mm]}]
  \def\ltbseg#1#2#3#4#5{%
    \draw[fill=#4] (#1,#3-2.5) rectangle (#2,#3+2.5);
    \node[lab] at ({(#1+#2)/2},#3) {#5};}
  \node[row] at (14,10) {CPU};
  \ltbseg{16}{44}{10}{white}{\iflangja{スレッド T}{thread T}}
  \ltbseg{44}{61}{10}{ln-accent!22}{\iflangja{トップハーフ}{top half}}
  \ltbseg{61}{80}{10}{white}{\iflangja{スレッド T}{thread T}}
  \ltbseg{80}{104}{10}{ln-box}{\iflangja{ボトムハーフ}{bottom half}}
  \ltbseg{104}{116}{10}{white}{T}
  % interrupt arrival
  \draw[->, thick, ln-caution] (44,20) -- (44,12.5);
  \node[lab, anchor=south, text=ln-caution] at (44,20) {\iflangja{割り込み}{interrupt}};
  % top half defers the bottom half
  \draw[->, ln-accent] (52.5,7.5) .. controls (56,1) and (86,1) .. (92,7.5);
  \node[lab, text=ln-accent] at (72,4.8) {\iflangja{後で実行するよう登録}{deferred}};
  % annotations
  \node[lab, anchor=north] at (52.5,-1)
    {\iflangja{直ちに実行\\高速・ブロックしない}{runs immediately;\\fast, non-blocking}};
  \node[lab, anchor=north] at (92,-1)
    {\iflangja{スレッドとしてスケジュール\\ブロック可能}{scheduled like a thread;\\may block}};
  \draw[->] (92,20) -- node[lab, above] {\iflangja{時間}{time}} (116,20);
\end{tikzpicture}
   (width ~116 mm)
\begin{tikzpicture}[x=1mm, y=1mm, font=\footnotesize\sffamily,
  lab/.style={font=\scriptsize\sffamily, align=center},
  row/.style={font=\scriptsize\sffamily, anchor=east},
  >={Stealth[length=1.6mm]}]
  \def\ltbseg#1#2#3#4#5{%
    \draw[fill=#4] (#1,#3-2.5) rectangle (#2,#3+2.5);
    \node[lab] at ({(#1+#2)/2},#3) {#5};}
  \node[row] at (14,10) {CPU};
  \ltbseg{16}{44}{10}{white}{\iflangja{スレッド T}{thread T}}
  \ltbseg{44}{61}{10}{ln-accent!22}{\iflangja{トップハーフ}{top half}}
  \ltbseg{61}{80}{10}{white}{\iflangja{スレッド T}{thread T}}
  \ltbseg{80}{104}{10}{ln-box}{\iflangja{ボトムハーフ}{bottom half}}
  \ltbseg{104}{116}{10}{white}{T}
  % interrupt arrival
  \draw[->, thick, ln-caution] (44,20) -- (44,12.5);
  \node[lab, anchor=south, text=ln-caution] at (44,20) {\iflangja{割り込み}{interrupt}};
  % top half defers the bottom half
  \draw[->, ln-accent] (52.5,7.5) .. controls (56,1) and (86,1) .. (92,7.5);
  \node[lab, text=ln-accent] at (72,4.8) {\iflangja{後で実行するよう登録}{deferred}};
  % annotations
  \node[lab, anchor=north] at (52.5,-1)
    {\iflangja{直ちに実行\\高速・ブロックしない}{runs immediately;\\fast, non-blocking}};
  \node[lab, anchor=north] at (92,-1)
    {\iflangja{スレッドとしてスケジュール\\ブロック可能}{scheduled like a thread;\\may block}};
  \draw[->] (92,20) -- node[lab, above] {\iflangja{時間}{time}} (116,20);
\end{tikzpicture}

  \caption{トップハーフとボトムハーフに分けた割り込み処理．トップハーフは
    スレッド T に割り込み，ボトムハーフはスケジューラが後で実行する．}
  \label{fig:l05-top-bottom}
\end{figure}

割り込みをスレッドとして処理することには代償がある．SunOS~5.0 では，
割り込みごとにおよそ 40 個の SPARC 命令が加わる．一方で他のところでは作業
を節約する．ミューテックスを必要とするハンドラがブロックできるようになった
ので，ミューテックスの操作は，ミューテックスを保持している間ハンドラが
実行されないようにするために割り込みマスクを変更する必要がなくなり，
ミューテックス操作1回当たりおよそ 12 命令が節約される．割り込み1回当たり
$c$ 命令が加わり，ミューテックス操作1回当たり $s$ 命令が節約されるとき，
この方式が引き合うのは
\begin{equation}
  c\, N_{\mathrm{int}} \;<\; s\, N_{\mathrm{mutex}}
  \quad\Longleftrightarrow\quad
  \frac{N_{\mathrm{int}}}{N_{\mathrm{mutex}}} \;<\; \frac{s}{c},
  \label{eq:l05-break-even}
\end{equation}
のときである．ここで $N_{\mathrm{int}}$ と $N_{\mathrm{mutex}}$ は，
割り込みとミューテックス操作の回数である．ミューテックスのロックとアン
ロックは割り込みの発生よりはるかに頻繁なので，収支は正になる．%
\margin{Linux も同じ取引をしている．\texttt{threadirqs} を指定して起動すると
ハンドラはスレッドで実行され，Linux~6.12 で本流に入った
\texttt{PREEMPT\_RT} では，\texttt{IRQF\_NO\_THREAD} の付かないハンドラは
すべてこれが既定になる．}

\begin{example}
  あるシステムが毎秒 2\,000 回の割り込みを処理し，150\,000 回の
  ミューテックス操作を行うとする．$c = 40$，$s = 12$ とすると，割り込み
  スレッドは毎秒 $40 \times 2\,000 = 80\,000$ 命令を費やし，
  $12 \times 150\,000 = 1\,800\,000$ 命令を節約するので，差し引き毎秒
  $1\,720\,000$ 命令の節約になる．\cref{eq:l05-break-even} により，この
  方式が損になるのは，ミューテックス操作当たりの割り込みが
  $12/40 = 0.3$ 回を超えるとき，すなわちこのミューテックス頻度では毎秒
  45\,000 回を超える割り込みがあるときだけである．
\end{example}

\begin{keypoint}
  頻出ケースを最適化せよ．スレッドとしての割り込みは，稀な事象である
  割り込みをわずかに高価にする代わりに，頻繁な事象であるミューテックス
  操作を安価にする．
\end{keypoint}

\section{スレッドとシグナル処理}
\label{sec:l05-thread-signals}

多対多の実装では，シグナルマスクは2つのレベルに存在する．\source{P2L4，
スレッドとシグナル処理，ケース~1--4；Stein と Shah \cite{stein1992}．}%
ユーザレベルライブラリはユーザレベルスレッドごとのマスクをそのスレッド
データ構造に保持し，カーネルはカーネルレベルスレッドごとのマスクをその LWP
に保持する．ユーザレベルスレッドは自分のマスクを，システムコールなしに
ライブラリを通じて変更する．カーネルにはカーネルレベルスレッドのマスクしか
見えない．シグナルが到着すると，カーネルは，その上で現在実行中のユーザ
レベルスレッドのマスクがどうであれ，マスクがそのシグナルを有効にしている
カーネルレベルスレッドに割り込む．\margin{NPTL では2つのレベルが一致するので，
この食い違いは生じない．\texttt{pthread\_sigmask} はシステムコール
（\texttt{rt\_sigprocmask}）であり，そのスレッドに対するカーネルのマスクを
直接設定する．}

SunOS のライブラリは，プロセスのシグナルハンドラテーブルに，それらの
シグナルのハンドラとして自分のルーチンを設定することで，この食い違いを
解消する．このルーチンが最初に実行され，ユーザレベルのマスクを調べて，
適切なユーザレベルスレッドでアプリケーションのハンドラを呼び出す．生じる
ケースは4つある（\cref{fig:l05-signal-cases}）．
\begin{description}
  \item[ケース 1] シグナルが，カーネルレベルスレッドでも，その上で実行中の
    ユーザレベルスレッドでも有効である．カーネルはカーネルレベルスレッドに
    割り込み，ライブラリのルーチンは現在のユーザレベルスレッドでシグナルが
    有効であることを見つけて，そこでハンドラを実行する．
  \item[ケース 2] 実行中のユーザレベルスレッドではシグナルが無効だが，
    実行可能で実行中ではない別のユーザレベルスレッドでは有効である．
    カーネルは，マスクがシグナルを有効にしているカーネルレベルスレッドに
    割り込む．ライブラリのルーチンは現在のスレッドでシグナルがマスクされて
    いることを知り，実行可能なそのスレッドをこのカーネルレベルスレッドに
    スケジュールする．ハンドラはその中で実行される．
  \item[ケース 3] シグナルが有効なユーザレベルスレッドが，別のカーネル
    レベルスレッド上で実行されている．割り込まれたカーネルレベルスレッド上
    のライブラリのルーチンは，そのスレッドがどこで実行されているかを知って
    いるので，もう一方のカーネルレベルスレッドを宛先に\emph{指定した}
    シグナルを送るシステムコールを行う．カーネルはそのカーネルレベル
    スレッドに割り込み，そこでライブラリのルーチンはシグナルが有効である
    ことを見つけてハンドラを実行する．
  \item[ケース 4] シグナルを有効にしているユーザレベルスレッドは1つもない
    が，カーネルレベルスレッドのマスクはまだそれを有効にしている．割り込ま
    れたカーネルレベルスレッド上のライブラリのルーチンは，そのカーネル
    レベルスレッドのマスクでシグナルを無効にするシステムコールを行い，
    プロセスに対してシグナルを再送出する．カーネルはそれを，まだ有効に
    している別のカーネルレベルスレッドに配送し，そこでも同じ手順が続く．
    これは，すべてのカーネルレベルスレッドでシグナルが無効になり，シグナル
    が保留のまま残るまで繰り返される．後にあるユーザレベルスレッドがその
    シグナルを有効にすると，ライブラリはシステムコールを行ってカーネル
    レベルスレッドでも再び有効にし，保留されていたシグナルが配送される．%
    \tip{4つのケースは1つの規則から生じる．ユーザレベルのマスクがシグナルを
    有効にしているところでハンドラを実行し，どのスレッドも有効にしていなけれ
    ば，カーネルレベルのマスクを落として再送出する．}
\end{description}

\begin{figure}[htbp]
  \centering
  % Signal handling with user-level (U) and kernel-level (K) signal masks:
% the four cases handled by the SunOS user-level thread library.
% Usage: % Signal handling with user-level (U) and kernel-level (K) signal masks:
% the four cases handled by the SunOS user-level thread library.
% Usage: % Signal handling with user-level (U) and kernel-level (K) signal masks:
% the four cases handled by the SunOS user-level thread library.
% Usage: \input{figures/l05-signal-cases}   (width ~116 mm)
\begin{tikzpicture}[x=1mm, y=1mm, font=\scriptsize\sffamily,
  >={Stealth[length=1.5mm]},
  ut/.style={draw, circle, fill=white, minimum size=5mm, inner sep=0pt},
  kt/.style={draw, rectangle, fill=ln-box, minimum size=5mm, inner sep=0pt},
  on/.style={font=\scriptsize\sffamily\bfseries, text=ln-tip, anchor=west},
  off/.style={font=\scriptsize\sffamily\bfseries, text=ln-caution, anchor=west},
  lab/.style={font=\scriptsize\sffamily, align=center},
  act/.style={->, thick, ln-accent},
  sig/.style={->, thick, ln-caution},
  ttl/.style={font=\scriptsize\sffamily, anchor=north west}]
  % \lsframe{title}: panel frame, title, user/kernel boundary, signal into K1
  \def\lsframe#1{%
    \draw[ln-rule, rounded corners=2pt] (0,-6) rectangle (56,31);
    \node[ttl] at (1,30.5) {#1};
    \draw[dashed, ln-rule] (1,12) -- (55,12);
    \draw[sig] (12,-5) -- (12,2.5);
    \node[lab, text=ln-caution, anchor=east] at (11,-2.5) {\iflangja{シグナル}{signal}};}
  % ---------------------------------------------------------------- case 1
  \begin{scope}[shift={(0,40)}]
    \lsframe{\textbf{\iflangja{ケース 1}{Case 1}}\enspace
      \iflangja{実行中の U1 が有効}{running U1 enabled}}
    \node[ut] (u1) at (12,19) {U1}; \node[on] at (14.8,19) {1};
    \node[kt] (k1) at (12,5) {K1};  \node[on] at (14.8,5) {1};
    \draw (u1) -- (k1);
    \node[lab, anchor=west, align=left] at (22,19)
      {\iflangja{ライブラリ：U1 は有効\\→ U1 でハンドラを実行}%
                {library: U1 enabled\\$\to$ handler runs in U1}};
  \end{scope}
  % ---------------------------------------------------------------- case 2
  \begin{scope}[shift={(60,40)}]
    \lsframe{\textbf{\iflangja{ケース 2}{Case 2}}\enspace
      \iflangja{実行可能な U2 が有効}{runnable U2 enabled}}
    \node[ut] (u1) at (12,19) {U1}; \node[off] at (14.8,19) {0};
    \node[kt] (k1) at (12,5) {K1};  \node[on] at (14.8,5) {1};
    \draw (u1) -- (k1);
    \node[ut, dashed] (u2) at (40,21) {U2}; \node[on] at (42.8,21) {1};
    \node[lab, text=ln-muted, anchor=north] at (40,18) {\iflangja{実行待ち}{ready}};
    \draw[act, dashed] (35,15.5) -- node[lab, below right, pos=0.5, text=ln-accent]
      {\iflangja{K1 に割り当て}{scheduled on K1}} (18.5,6.5);
  \end{scope}
  % ---------------------------------------------------------------- case 3
  \begin{scope}[shift={(0,0)}]
    \lsframe{\textbf{\iflangja{ケース 3}{Case 3}}\enspace
      \iflangja{別の K 上の U2 が有効}{U2 on another K enabled}}
    \node[ut] (u1) at (12,19) {U1}; \node[off] at (14.8,19) {0};
    \node[kt] (k1) at (12,5) {K1};  \node[on] at (14.8,5) {1};
    \draw (u1) -- (k1);
    \node[ut] (u2) at (42,19) {U2}; \node[on] at (44.8,19) {1};
    \node[kt] (k2) at (42,5) {K2};  \node[on] at (44.8,5) {1};
    \draw (u2) -- (k2);
    \draw[act] (19,5) -- node[lab, below, text=ln-accent]
      {\iflangja{指定した\\シグナル}{directed\\signal}} (39.2,5);
  \end{scope}
  % ---------------------------------------------------------------- case 4
  \begin{scope}[shift={(60,0)}]
    \lsframe{\textbf{\iflangja{ケース 4}{Case 4}}\enspace
      \iflangja{有効な U がない}{no U enabled}}
    \node[ut] (u1) at (12,19) {U1}; \node[off] at (14.8,19) {0};
    \node[kt] (k1) at (12,5) {K1};
    \node[font=\scriptsize\sffamily\bfseries, anchor=west] at (14.8,5)
      {\textcolor{ln-tip}{1}$\to$\textcolor{ln-caution}{0}};
    \draw (u1) -- (k1);
    \node[ut] (u2) at (42,19) {U2}; \node[off] at (44.8,19) {0};
    \node[kt] (k2) at (42,5) {K2};
    \node[font=\scriptsize\sffamily\bfseries, anchor=west] at (44.8,5)
      {\textcolor{ln-tip}{1}$\to$\textcolor{ln-caution}{0}};
    \draw (u2) -- (k2);
    \draw[act] (24,5) -- node[lab, below, text=ln-accent]
      {\iflangja{再送出}{re-raised}} (39.2,5);
  \end{scope}
\end{tikzpicture}
   (width ~116 mm)
\begin{tikzpicture}[x=1mm, y=1mm, font=\scriptsize\sffamily,
  >={Stealth[length=1.5mm]},
  ut/.style={draw, circle, fill=white, minimum size=5mm, inner sep=0pt},
  kt/.style={draw, rectangle, fill=ln-box, minimum size=5mm, inner sep=0pt},
  on/.style={font=\scriptsize\sffamily\bfseries, text=ln-tip, anchor=west},
  off/.style={font=\scriptsize\sffamily\bfseries, text=ln-caution, anchor=west},
  lab/.style={font=\scriptsize\sffamily, align=center},
  act/.style={->, thick, ln-accent},
  sig/.style={->, thick, ln-caution},
  ttl/.style={font=\scriptsize\sffamily, anchor=north west}]
  % \lsframe{title}: panel frame, title, user/kernel boundary, signal into K1
  \def\lsframe#1{%
    \draw[ln-rule, rounded corners=2pt] (0,-6) rectangle (56,31);
    \node[ttl] at (1,30.5) {#1};
    \draw[dashed, ln-rule] (1,12) -- (55,12);
    \draw[sig] (12,-5) -- (12,2.5);
    \node[lab, text=ln-caution, anchor=east] at (11,-2.5) {\iflangja{シグナル}{signal}};}
  % ---------------------------------------------------------------- case 1
  \begin{scope}[shift={(0,40)}]
    \lsframe{\textbf{\iflangja{ケース 1}{Case 1}}\enspace
      \iflangja{実行中の U1 が有効}{running U1 enabled}}
    \node[ut] (u1) at (12,19) {U1}; \node[on] at (14.8,19) {1};
    \node[kt] (k1) at (12,5) {K1};  \node[on] at (14.8,5) {1};
    \draw (u1) -- (k1);
    \node[lab, anchor=west, align=left] at (22,19)
      {\iflangja{ライブラリ：U1 は有効\\→ U1 でハンドラを実行}%
                {library: U1 enabled\\$\to$ handler runs in U1}};
  \end{scope}
  % ---------------------------------------------------------------- case 2
  \begin{scope}[shift={(60,40)}]
    \lsframe{\textbf{\iflangja{ケース 2}{Case 2}}\enspace
      \iflangja{実行可能な U2 が有効}{runnable U2 enabled}}
    \node[ut] (u1) at (12,19) {U1}; \node[off] at (14.8,19) {0};
    \node[kt] (k1) at (12,5) {K1};  \node[on] at (14.8,5) {1};
    \draw (u1) -- (k1);
    \node[ut, dashed] (u2) at (40,21) {U2}; \node[on] at (42.8,21) {1};
    \node[lab, text=ln-muted, anchor=north] at (40,18) {\iflangja{実行待ち}{ready}};
    \draw[act, dashed] (35,15.5) -- node[lab, below right, pos=0.5, text=ln-accent]
      {\iflangja{K1 に割り当て}{scheduled on K1}} (18.5,6.5);
  \end{scope}
  % ---------------------------------------------------------------- case 3
  \begin{scope}[shift={(0,0)}]
    \lsframe{\textbf{\iflangja{ケース 3}{Case 3}}\enspace
      \iflangja{別の K 上の U2 が有効}{U2 on another K enabled}}
    \node[ut] (u1) at (12,19) {U1}; \node[off] at (14.8,19) {0};
    \node[kt] (k1) at (12,5) {K1};  \node[on] at (14.8,5) {1};
    \draw (u1) -- (k1);
    \node[ut] (u2) at (42,19) {U2}; \node[on] at (44.8,19) {1};
    \node[kt] (k2) at (42,5) {K2};  \node[on] at (44.8,5) {1};
    \draw (u2) -- (k2);
    \draw[act] (19,5) -- node[lab, below, text=ln-accent]
      {\iflangja{指定した\\シグナル}{directed\\signal}} (39.2,5);
  \end{scope}
  % ---------------------------------------------------------------- case 4
  \begin{scope}[shift={(60,0)}]
    \lsframe{\textbf{\iflangja{ケース 4}{Case 4}}\enspace
      \iflangja{有効な U がない}{no U enabled}}
    \node[ut] (u1) at (12,19) {U1}; \node[off] at (14.8,19) {0};
    \node[kt] (k1) at (12,5) {K1};
    \node[font=\scriptsize\sffamily\bfseries, anchor=west] at (14.8,5)
      {\textcolor{ln-tip}{1}$\to$\textcolor{ln-caution}{0}};
    \draw (u1) -- (k1);
    \node[ut] (u2) at (42,19) {U2}; \node[off] at (44.8,19) {0};
    \node[kt] (k2) at (42,5) {K2};
    \node[font=\scriptsize\sffamily\bfseries, anchor=west] at (44.8,5)
      {\textcolor{ln-tip}{1}$\to$\textcolor{ln-caution}{0}};
    \draw (u2) -- (k2);
    \draw[act] (24,5) -- node[lab, below, text=ln-accent]
      {\iflangja{再送出}{re-raised}} (39.2,5);
  \end{scope}
\end{tikzpicture}
   (width ~116 mm)
\begin{tikzpicture}[x=1mm, y=1mm, font=\scriptsize\sffamily,
  >={Stealth[length=1.5mm]},
  ut/.style={draw, circle, fill=white, minimum size=5mm, inner sep=0pt},
  kt/.style={draw, rectangle, fill=ln-box, minimum size=5mm, inner sep=0pt},
  on/.style={font=\scriptsize\sffamily\bfseries, text=ln-tip, anchor=west},
  off/.style={font=\scriptsize\sffamily\bfseries, text=ln-caution, anchor=west},
  lab/.style={font=\scriptsize\sffamily, align=center},
  act/.style={->, thick, ln-accent},
  sig/.style={->, thick, ln-caution},
  ttl/.style={font=\scriptsize\sffamily, anchor=north west}]
  % \lsframe{title}: panel frame, title, user/kernel boundary, signal into K1
  \def\lsframe#1{%
    \draw[ln-rule, rounded corners=2pt] (0,-6) rectangle (56,31);
    \node[ttl] at (1,30.5) {#1};
    \draw[dashed, ln-rule] (1,12) -- (55,12);
    \draw[sig] (12,-5) -- (12,2.5);
    \node[lab, text=ln-caution, anchor=east] at (11,-2.5) {\iflangja{シグナル}{signal}};}
  % ---------------------------------------------------------------- case 1
  \begin{scope}[shift={(0,40)}]
    \lsframe{\textbf{\iflangja{ケース 1}{Case 1}}\enspace
      \iflangja{実行中の U1 が有効}{running U1 enabled}}
    \node[ut] (u1) at (12,19) {U1}; \node[on] at (14.8,19) {1};
    \node[kt] (k1) at (12,5) {K1};  \node[on] at (14.8,5) {1};
    \draw (u1) -- (k1);
    \node[lab, anchor=west, align=left] at (22,19)
      {\iflangja{ライブラリ：U1 は有効\\→ U1 でハンドラを実行}%
                {library: U1 enabled\\$\to$ handler runs in U1}};
  \end{scope}
  % ---------------------------------------------------------------- case 2
  \begin{scope}[shift={(60,40)}]
    \lsframe{\textbf{\iflangja{ケース 2}{Case 2}}\enspace
      \iflangja{実行可能な U2 が有効}{runnable U2 enabled}}
    \node[ut] (u1) at (12,19) {U1}; \node[off] at (14.8,19) {0};
    \node[kt] (k1) at (12,5) {K1};  \node[on] at (14.8,5) {1};
    \draw (u1) -- (k1);
    \node[ut, dashed] (u2) at (40,21) {U2}; \node[on] at (42.8,21) {1};
    \node[lab, text=ln-muted, anchor=north] at (40,18) {\iflangja{実行待ち}{ready}};
    \draw[act, dashed] (35,15.5) -- node[lab, below right, pos=0.5, text=ln-accent]
      {\iflangja{K1 に割り当て}{scheduled on K1}} (18.5,6.5);
  \end{scope}
  % ---------------------------------------------------------------- case 3
  \begin{scope}[shift={(0,0)}]
    \lsframe{\textbf{\iflangja{ケース 3}{Case 3}}\enspace
      \iflangja{別の K 上の U2 が有効}{U2 on another K enabled}}
    \node[ut] (u1) at (12,19) {U1}; \node[off] at (14.8,19) {0};
    \node[kt] (k1) at (12,5) {K1};  \node[on] at (14.8,5) {1};
    \draw (u1) -- (k1);
    \node[ut] (u2) at (42,19) {U2}; \node[on] at (44.8,19) {1};
    \node[kt] (k2) at (42,5) {K2};  \node[on] at (44.8,5) {1};
    \draw (u2) -- (k2);
    \draw[act] (19,5) -- node[lab, below, text=ln-accent]
      {\iflangja{指定した\\シグナル}{directed\\signal}} (39.2,5);
  \end{scope}
  % ---------------------------------------------------------------- case 4
  \begin{scope}[shift={(60,0)}]
    \lsframe{\textbf{\iflangja{ケース 4}{Case 4}}\enspace
      \iflangja{有効な U がない}{no U enabled}}
    \node[ut] (u1) at (12,19) {U1}; \node[off] at (14.8,19) {0};
    \node[kt] (k1) at (12,5) {K1};
    \node[font=\scriptsize\sffamily\bfseries, anchor=west] at (14.8,5)
      {\textcolor{ln-tip}{1}$\to$\textcolor{ln-caution}{0}};
    \draw (u1) -- (k1);
    \node[ut] (u2) at (42,19) {U2}; \node[off] at (44.8,19) {0};
    \node[kt] (k2) at (42,5) {K2};
    \node[font=\scriptsize\sffamily\bfseries, anchor=west] at (44.8,5)
      {\textcolor{ln-tip}{1}$\to$\textcolor{ln-caution}{0}};
    \draw (u2) -- (k2);
    \draw[act] (24,5) -- node[lab, below, text=ln-accent]
      {\iflangja{再送出}{re-raised}} (39.2,5);
  \end{scope}
\end{tikzpicture}

  \caption{ユーザレベルスレッド（U）とカーネルレベルスレッド（K）を用いた
    シグナル処理の4つのケース．各スレッドの脇の数字は，そのシグナルに
    対するマスクのビットであり，1 が有効，0 が無効を表す．}
  \label{fig:l05-signal-cases}
\end{figure}

この設計は，スレッドとしての割り込みと同じ原理に従う．シグナルの発生は，
スレッドがシグナルマスクを変更する頻度---典型的にはクリティカルセクション
ごとである---よりはるかに低い．ユーザレベルのマスクの変更は安価なライブラリ
操作であるのに対し，カーネルレベルのマスクをそれに合わせて保つには，変更の
たびにシステムコールが必要になる．その代わりにシグナルの配送は高価になる
---ライブラリのルーチンが最初に実行され，シグナルを転送したりシステム
コールを行ったりしなければならないことがある---が，これは稀にしか起こら
ない．

\section{Linux におけるタスク}
\label{sec:l05-linux}

Linux カーネルはプロセスとスレッドを区別しない．\source{P2L4，Linux に
おけるタスク；Silberschatz ら ch.~4；Kerrisk ch.~28．}どちらも1種類の実行
コンテキストから作られ，マルチスレッドプロセスは資源を共有するそうした
コンテキストの集まりである．

\begin{definition}[タスク]
  Linux カーネルにおける\term[たすく]{タスク}[task]とは，基本的な実行
  コンテキストであり，\texttt{struct task\_struct} で表されるカーネル
  レベルスレッドである．シングルスレッドプロセスは1つのタスクからなり，
  マルチスレッドプロセスは複数のタスクからなる．
\end{definition}

\needspace{13\baselineskip}
この構造体の主なフィールドは次のとおりである．
\begin{lstlisting}[language=C, numbers=none]
struct task_struct {       // 抜粋，Linux 3.x
    // ...
    pid_t pid;             // タスクの識別子
    pid_t tgid;            // スレッドグループの識別子
    int prio;
    volatile long state;
    struct mm_struct *mm;         // アドレス空間
    struct files_struct *files;   // 開いたファイル
    struct list_head tasks;
    int on_cpu;
    cpumask_t cpus_allowed;
    // ...
};
\end{lstlisting}
すべてのタスクは独自の識別子 \texttt{pid} を持つ．1つのプロセスのタスクは
スレッドグループを形成し，その識別子 \texttt{tgid} はグループの先頭タスク
の \texttt{pid} である．これが \texttt{getpid} の返すプロセス識別子である．%
\caution{\texttt{pid} という名前のフィールドはスレッドの識別子である．
\texttt{getpid} は \texttt{tgid} を返し，\texttt{gettid} が \texttt{pid} を
返す．\texttt{ps -L} はこの2つを \texttt{PID} と \texttt{LWP} として表示
する．}
タスク構造体はプロセスの資源そのものを含まない．\texttt{mm} と
\texttt{files} は，アドレス空間と開いたファイルのために別に割り当てられた
構造体を指す．\cref{sec:l05-structures}と同様に，状態は複数のタスクが共有
できる構造に分けられている．

タスクは \term{clone}[clone] システムコールで生成される．
\begin{lstlisting}[language=C, numbers=none]
clone(function, stack_ptr, sharing_flags, args);
\end{lstlisting}
新しいタスクはスタック \texttt{stack\_ptr} 上で \texttt{function(args)} を
実行する．フラグ \texttt{sharing\_flags} の各ビットは，状態の一部を生成元の
タスクと共有するか複製するかを選ぶ（\cref{tab:l05-clone}）．これらの
フラグをすべて立て，さらに新しいタスクを生成元のスレッドグループに入れる
\texttt{CLONE\_THREAD} を立てると，新しいタスクは生成したプロセスの
スレッドになる．すべてを落とすと新しいプロセスになり，\texttt{fork} は
このように実装されている．

\begin{table}[htbp]
  \centering
  \caption{\texttt{clone} の共有フラグ．}
  \label{tab:l05-clone}
  \small
  \begin{tabular}{@{}l>{\raggedright\arraybackslash}p{0.33\linewidth}>{\raggedright\arraybackslash}p{0.33\linewidth}@{}}
    \toprule
    フラグ & 立てた場合 & 落とした場合 \\
    \midrule
    \texttt{CLONE\_FILES}   & ファイルディスクリプタ表を共有
                            & ファイルディスクリプタ表を複製 \\
    \texttt{CLONE\_FS}      & ファイルシステム情報を共有
                            & ファイルシステム情報を複製 \\
    \texttt{CLONE\_SIGHAND} & シグナルハンドラテーブルを共有
                            & シグナルハンドラテーブルを複製 \\
    \texttt{CLONE\_VM}      & アドレス空間を共有
                            & 新しいアドレス空間（複製） \\
    \bottomrule
  \end{tabular}
\end{table}

現在の Linux システムにおける POSIX スレッドは，
\term{NPTL}[Native POSIX Thread Library]%
\index{Native POSIX Thread Library|see{NPTL}}（Native POSIX Thread
Library）によって実装されており，NPTL は一対一モデルを用いる．すなわち，
各ユーザレベルスレッドが1つのタスクである．したがってカーネルはすべての
スレッドを，そのシグナルマスクなどの情報とともに見ることができ，
\cref{sec:l05-interaction}と\cref{sec:l05-thread-signals}の協調の問題は
生じない．かつて一対一モデルに反対する論拠となっていた欠点は薄れた．カーネルへの
トラップははるかに安価になり，メモリとスレッド識別子の範囲は多数の
カーネルレベルスレッドに十分なほど大きい．NPTL は LinuxThreads を置き
換えた．LinuxThreads も1スレッドにつき1タスクを生成していたが，シグナル
処理などの点で POSIX に準拠していなかった．同じ時期に開発された多対多の
ライブラリである Next Generation POSIX Threads（NGPT）は
\cref{sec:l05-interaction}の協調の問題を抱えており，NPTL のために放棄
された．\margin{LinuxThreads は多対多と説明されることがあるが，当時の
多対多ライブラリは NGPT である．}

\begin{keypoint}[章のまとめ]
  スレッドの状態は，すべてのスレッドが共有するハードプロセス状態，
  カーネルレベルスレッドごとのライトプロセス状態，そしてレッドゾーンを
  伴うユーザレベルの構造に分けられる．SunOS~5.0 では LWP が2つのレベルを
  つなぎ，両者はシステムコールとシグナルで協調する．ハンドラは割り込まれた
  スレッドの文脈で実行されるので，マスクがデッドロックを防ぐ．スレッドと
  しての割り込み，トップハーフとボトムハーフ，ユーザレベルのシグナルマスク
  は，いずれも頻出ケースを最適化する．Linux はすべてのスレッドを
  \texttt{clone} で生成されるタスクとし，NPTL がそれを一対一に対応付ける．
\end{keypoint}

\end{document}
